\chapter{多重积分与积分变换}\label{ch:c4}

细丝的质量、沿曲线的功和穿过曲面的通量，虽然都由微元累积得到，却使用不同的长度、方向和权重。换用积分变量时，区域边界与微元一同变换。使用积分定理时，还须给出边界方向，并确认场在相关区域内光滑。

本讲从线积分和曲面积分进入多重积分，再利用高斯积分引出Gamma与Beta函数。傅里叶展开把信号分解为频率分量，拉普拉斯变换把初值和微分方程放入同一个代数式。傅里叶正反变换采用配套的归一化系数，拉普拉斯变换采用单边初值约定；谐波、延迟输入和卷积均按这些约定计算。


\section{曲线积分与曲面积分}\label{sec:c4-line}
沿细线求质量、沿路径求功和穿过曲面求通量，使用的有向性不同。先由所求物理量确定积分元，再参数化曲线或曲面，方向约定随参数一并说明。

\subsection{弧长权重与有向功}
同一条曲线可以承载两类不同积分。标量弧长积分$\int_Cf\dd s$把线密度等标量乘以弧长微元相加，不随曲线方向改变；向量场的功$\int_C\vct F\cdot\dd\vct r$则使用有向位移，反向遍历会改变符号。

取参数曲线$\vct r(t)$，$a\le t\le b$，有
\[\int_C f\dd s=\int_a^b f(\vct r(t))|\vct r'(t)|\dd t,\qquad
\int_C\vct F\cdot\dd\vct r=\int_a^b\vct F(\vct r(t))\cdot\vct r'(t)\dd t.
\]
这里$a<b$且曲线方向由$t$增加确定。弧长的模长不能漏掉；做功积分则不能把向量线元一律替换成正的弧长。

直角坐标中$\dd\vct r=\vct e_x\dd x+\vct e_y\dd y+\vct e_z\dd z$。柱坐标中
\[\dd\vct r=\vct e_r\dd r+\vct e_\varphi r\dd\varphi+\vct e_z\dd z,\quad
\dd s^2=\dd r^2+r^2\dd\varphi^2+\dd z^2.
\]
球坐标取极角$\theta$、方位角$\varphi$，则
\[\dd\vct r=\vct e_r\dd r+\vct e_\theta r\dd\theta+\vct e_\varphi r\sin\theta\dd\varphi,
\quad\dd s^2=\dd r^2+r^2\dd\theta^2+r^2\sin^2\theta\dd\varphi^2.
\]
参数化可以更换，实际积分不变；是否简化应看曲线约束和被积函数，而非只看坐标名称。

\Needspace{8\baselineskip}
\subsection{曲线参数与积分方向}
把曲线分成许多小弧段，标量线积分用每段上的密度乘实际弧长。若曲线参数为$t$，小弧长为$|\vct r'(t)|\dd t$，这个模长把参数宽度转成空间长度。倒向走过同一段细线不会改变它的质量，所以标量弧长积分与方向无关。

功积分则把每一微小位移与力作点积。沿曲线单位切向记为$\vct t$，便有$\vct F\cdot\dd\vct r=(\vct F\cdot\vct t)\dd s$，只有切向力分量做功。反向时单位切向反转，所以相同场沿同一几何曲线的功变号。标量弧长积分使用正长度，功积分使用有向位移，反向遍历时因而具有不同符号。

重新参数化不改变积分值。若$t=t(u)$单调增加，速度链式法则使$|\dd\vct r/\dd u|=|\vct r'(t)|t'(u)$，正好与$\dd t=t'(u)\dd u$配对。若参数倒向，弧长仍取绝对值，有向功却保留负号。参数停留或重复遍历时，还须说明是否在计算真实材料总量，还是沿运动路径累计了多次。

柱、球坐标中的位移微元应使用单位基矢的物理分量。固定$r$改变$\varphi$，移动长度是$r\dd\varphi$；固定$r,\theta$改变球坐标方位角，长度则为$r\sin\theta\dd\varphi$。坐标函数本身的导数不总等于实际方向速度，尺度因子需与场分量一起进入点积。例~\ref{ex:c4-e1}沿一条射线，角微分均为零，正适合比较三种坐标表达。

\begin{example}[同一线积分的三种坐标]\label{ex:c4-e1}
线段$C$从$(0,0,0)$到$(1,2,1)$，求$\int_C(x^2-y+3z)\dd s$。分别用直角、柱、球坐标说明线元。坐标数值按同一长度单位表示。 见图~\ref{fig:c4-e1-q}。
\end{example}

\begin{figure}[H]\centering
\begin{tikzpicture}[>=Stealth,line cap=round,line join=round]
\draw[axis](0,0)--(3,0)node[right]{$x$};\draw[axis](0,0)--(-1.2,-.8)node[left]{$y$};\draw[axis](0,0)--(0,2.6)node[above]{$z$};\draw[axis](0,0)--(1.4,1.4)node[above right]{$B(1,2,1)$};\node[below]at(0,0){$O$};
\end{tikzpicture}
\caption{原点出发的直线方向固定，三种坐标都可用单一径向参数。}\label{fig:c4-e1-q}
\end{figure}
\begin{solution}
沿图~\ref{fig:c4-e1-q}的线段取
$\vct r(u)=(u,2u,u)$，$0\le u\le1$。
在直角坐标中，$\dd\vct r=(1,2,1)\dd u$，所以
$\dd s=|\dd\vct r|=\sqrt6\dd u$。
另外两种坐标描述同一条射线：
\[
 \begin{aligned}
 \text{柱坐标：}\quad&
 \varrho=\sqrt5u,\quad\varphi=\arctan(2),\quad z=u,\\
 &\dd s=\sqrt{\dd\varrho^2+\varrho^2\dd\varphi^2+\dd z^2}
 =\sqrt{5+1}\dd u;\\
 \text{球坐标：}\quad&
 R=\sqrt6u,\quad\theta=\arccos(1/\sqrt6),\quad
 \varphi=\arctan(2),\\
 &\dd s=\sqrt{\dd R^2+R^2\dd\theta^2+
 R^2\sin^2\theta\dd\varphi^2}=\dd R.
 \end{aligned}
\]
两种角坐标沿开线段均不变。原点处角度虽不唯一，可按该射线取极限，不影响积分。
\begin{figure}[H]\centering
\begin{tikzpicture}[>=Stealth,line cap=round,line join=round]
\draw[axis](0,0)--(3.5,0)node[right]{$u$};\draw[axis](0,0)--(0,2.5)node[above]{$s$};\draw[curve](0,0)--(1,2.44949);\draw[guide](.5,1.22474)--(.8,1.22474)--(.8,1.95959);\node[below]at(.65,1.22474){$\dd u$};\node[right]at(.8,1.6){$\sqrt6\dd u$};
\end{tikzpicture}
\caption{线段的弧长与参数之比恒为$\sqrt6$。}\label{fig:c4-e1-s}
\end{figure}
图~\ref{fig:c4-e1-s}标出的弧长增量与参数增量之比恒为$\sqrt6$。被积函数沿线化为
$x^2-y+3z=u^2+u$，故
\[
 \boxed{I=\sqrt6\int_0^1(u^2+u)\dd u
 =\sqrt6\left[\frac{u^3}{3}+\frac{u^2}{2}\right]_0^1
 =\frac{5\sqrt6}{6}}.
\]
作为定义上的比较，若求向量场$\vct A=(x^2,-y,3z)$沿这条线段的功，须用有向位移作点积：
\[
 W=\int_C\vct A\cdot\dd\vct r
 =\int_0^1(u^2-4u+3u)\dd u
 =\left[\frac{u^3}{3}-\frac{u^2}{2}\right]_0^1=-\frac16.
\]
反向沿线段运动时，$W$变号，标量弧长积分$I$不变。
\end{solution}

\Needspace{8\baselineskip}
\begin{example}[变密度螺旋线的质心]\label{ex:c4-e2}
细线参数为$\vct r(t)=\frac L{\sqrt2}(\cos t,\sin t,t)$，$0\le t\le6\pi$，线密度$\lambda=\lambda_0(1+z/L)$。求总质量与三个质心坐标。 见图~\ref{fig:c4-e2-q}。
\end{example}

\begin{figure}[H]\centering
\begin{tikzpicture}[>=Stealth,line cap=round,line join=round]
\draw[curve]plot coordinates{(1.0000000,0.0000000) (1.0055666,0.0339147) (1.0071943,0.0677427) (1.0049067,0.1013976) (0.9987425,0.1347937) (0.9887556,0.1678462) (0.9750154,0.2004719) (0.9576058,0.2325891) (0.9366253,0.2641182) (0.9121863,0.2949818) (0.8844152,0.3251054) (0.8534513,0.3544170) (0.8194465,0.3828481) (0.7825648,0.4103337) (0.7429815,0.4368124) (0.7008826,0.4622267) (0.6564640,0.4865236) (0.6099307,0.5096542) (0.5614961,0.5315744) (0.5113812,0.5522448) (0.4598134,0.5716308) (0.4070262,0.5897032) (0.3532574,0.6064377) (0.2987491,0.6218153) (0.2437463,0.6358226) (0.1884956,0.6484513) (0.1332449,0.6596987) (0.0782420,0.6695676) (0.0237337,0.6780660) (-0.0300350,0.6852076) (-0.0828223,0.6910113) (-0.1343901,0.6955014) (-0.1845050,0.6987071) (-0.2329395,0.7006631) (-0.2794728,0.7014085) (-0.3238915,0.7009878) (-0.3659904,0.6994495) (-0.4055737,0.6968469) (-0.4424554,0.6932375) (-0.4764602,0.6886824) (-0.5074241,0.6832469) (-0.5351952,0.6769995) (-0.5596341,0.6700120) (-0.5806147,0.6623590) (-0.5980243,0.6541179) (-0.6117645,0.6453683) (-0.6217513,0.6361919) (-0.6279156,0.6266719) (-0.6302032,0.6168931) (-0.6285754,0.6069412) (-0.6230089,0.5969026) (-0.6134958,0.5868640) (-0.6000439,0.5769121) (-0.5826767,0.5671333) (-0.5614328,0.5576134) (-0.5363663,0.5484369) (-0.5075464,0.5396873) (-0.4750572,0.5314462) (-0.4389970,0.5237932) (-0.3994784,0.5168057) (-0.3566277,0.5105583) (-0.3105841,0.5051228) (-0.2614996,0.5005677) (-0.2095383,0.4969583) (-0.1548754,0.4943557) (-0.0976968,0.4928174) (-0.0381985,0.4923967) (0.0234144,0.4931422) (0.0869286,0.4950981) (0.1521232,0.4983038) (0.2187706,0.5027939) (0.2866375,0.5085976) (0.3554859,0.5157392) (0.4250738,0.5242377) (0.4951563,0.5341065) (0.5654867,0.5453539) (0.6358170,0.5579826) (0.7058996,0.5719899) (0.7754875,0.5873675) (0.8443359,0.6041020) (0.9122028,0.6221744) (0.9788502,0.6415604) (1.0440447,0.6622308) (1.1075589,0.6841510) (1.1691719,0.7072816) (1.2286702,0.7315785) (1.2858487,0.7569928) (1.3405117,0.7834715) (1.3924730,0.8109571) (1.4415574,0.8393882) (1.4876010,0.8686998) (1.5304518,0.8988234) (1.5699703,0.9296870) (1.6060305,0.9612161) (1.6385198,0.9933333) (1.6673396,1.0259590) (1.6924061,1.0590115) (1.7136500,1.0924076) (1.7310173,1.1260625) (1.7444691,1.1598905) (1.7539822,1.1938052) (1.7595488,1.2277199) (1.7611766,1.2615479) (1.7588890,1.2952028) (1.7527247,1.3285989) (1.7427379,1.3616514) (1.7289977,1.3942771) (1.7115880,1.4263943) (1.6906075,1.4579234) (1.6661686,1.4887871) (1.6383975,1.5189106) (1.6074335,1.5482222) (1.5734287,1.5766533) (1.5365470,1.6041389) (1.4969637,1.6306176) (1.4548648,1.6560319) (1.4104462,1.6803288) (1.3639129,1.7034594) (1.3154783,1.7253796) (1.2653634,1.7460500) (1.2137957,1.7654360) (1.1610084,1.7835084) (1.1072396,1.8002429) (1.0527314,1.8156206) (0.9977285,1.8296278) (0.9424778,1.8422565) (0.8872271,1.8535039) (0.8322242,1.8633728) (0.7777159,1.8718712) (0.7239472,1.8790128) (0.6711599,1.8848166) (0.6195922,1.8893066) (0.5694773,1.8925123) (0.5210427,1.8944683) (0.4745094,1.8952138) (0.4300908,1.8947930) (0.3879919,1.8932547) (0.3484086,1.8906522) (0.3115269,1.8870427) (0.2775221,1.8824876) (0.2465581,1.8770521) (0.2187870,1.8708047) (0.1943481,1.8638172) (0.1733675,1.8561642) (0.1559579,1.8479231) (0.1422177,1.8391735) (0.1322309,1.8299971) (0.1260666,1.8204771) (0.1237790,1.8106983) (0.1254068,1.8007464) (0.1309734,1.7907078) (0.1404864,1.7806692) (0.1539383,1.7707173) (0.1713056,1.7609385) (0.1925495,1.7514186) (0.2176160,1.7422421) (0.2464358,1.7334925) (0.2789251,1.7252514) (0.3149853,1.7175984) (0.3545038,1.7106109) (0.3973546,1.7043635) (0.4433982,1.6989280) (0.4924826,1.6943730) (0.5444439,1.6907635) (0.5991069,1.6881609) (0.6562854,1.6866226) (0.7157837,1.6862019) (0.7773967,1.6869474) (0.8409109,1.6889033) (0.9061054,1.6921090) (0.9727528,1.6965991) (1.0406197,1.7024028) (1.1094681,1.7095444) (1.1790560,1.7180429) (1.2491386,1.7279117) (1.3194689,1.7391591) (1.3897993,1.7517878) (1.4598818,1.7657951) (1.5294697,1.7811727) (1.5983181,1.7979072) (1.6661850,1.8159796) (1.7328324,1.8353657) (1.7980270,1.8560360) (1.8615412,1.8779562) (1.9231541,1.9010868) (1.9826524,1.9253837) (2.0398310,1.9507981) (2.0944939,1.9772767) (2.1464552,2.0047623) (2.1955397,2.0331934) (2.2415832,2.0625051) (2.2844340,2.0926286) (2.3239526,2.1234922) (2.3600128,2.1550213) (2.3925020,2.1871385) (2.4213219,2.2197642) (2.4463883,2.2528167) (2.4676323,2.2862128) (2.4849995,2.3198677) (2.4984514,2.3536957) (2.5079645,2.3876104) (2.5135310,2.4215252) (2.5151588,2.4553532) (2.5128712,2.4890080) (2.5067069,2.5224041) (2.4967201,2.5554566) (2.4829799,2.5880823) (2.4655703,2.6201995) (2.4445897,2.6517286) (2.4201508,2.6825923) (2.3923797,2.7127158) (2.3614158,2.7420274) (2.3274110,2.7704585) (2.2905293,2.7979441) (2.2509460,2.8244228) (2.2088471,2.8498371) (2.1644284,2.8741340) (2.1178951,2.8972646) (2.0694606,2.9191848) (2.0193457,2.9398552) (1.9677779,2.9592412) (1.9149906,2.9773136) (1.8612219,2.9940481) (1.8067136,3.0094258) (1.7517107,3.0234330) (1.6964600,3.0360617) (1.6412093,3.0473091) (1.5862064,3.0571780) (1.5316982,3.0656764) (1.4779294,3.0728180) (1.4251422,3.0786218) (1.3735744,3.0831118) (1.3234595,3.0863176) (1.2750249,3.0882735) (1.2284916,3.0890190) (1.1840730,3.0885982) (1.1419741,3.0870599) (1.1023908,3.0844574) (1.0655091,3.0808479) (1.0315043,3.0762928) (1.0005404,3.0708573) (0.9727693,3.0646099) (0.9483303,3.0576224) (0.9273498,3.0499694) (0.9099402,3.0417283) (0.8962000,3.0329787) (0.8862131,3.0238023) (0.8800489,3.0142823) (0.8777612,3.0045035) (0.8793890,2.9945517) (0.8849556,2.9845130) (0.8944687,2.9744744) (0.9079205,2.9645225) (0.9252878,2.9547437) (0.9465317,2.9452238) (0.9715982,2.9360473) (1.0004180,2.9272977) (1.0329073,2.9190566) (1.0689675,2.9114037) (1.1084861,2.9044161) (1.1513368,2.8981687) (1.1973804,2.8927332) (1.2464648,2.8881782) (1.2984262,2.8845687) (1.3530891,2.8819661) (1.4102677,2.8804279) (1.4697660,2.8800071) (1.5313789,2.8807526) (1.5948931,2.8827085) (1.6600877,2.8859142) (1.7267350,2.8904043) (1.7946020,2.8962080) (1.8634504,2.9033496) (1.9330383,2.9118481) (2.0031208,2.9217169) (2.0734512,2.9329643) (2.1437815,2.9455930) (2.2138640,2.9596003) (2.2834519,2.9749779) (2.3523003,2.9917124) (2.4201673,3.0097848) (2.4868146,3.0291709) (2.5520092,3.0498412) (2.6155234,3.0717614) (2.6771363,3.0948920) (2.7366346,3.1191889) (2.7938132,3.1446033) (2.8484761,3.1710819) (2.9004375,3.1985675) (2.9495219,3.2269987) (2.9955655,3.2563103) (3.0384162,3.2864338) (3.0779348,3.3172974) (3.1139950,3.3488265) (3.1464843,3.3809437) (3.1753041,3.4135694) (3.2003706,3.4466220) (3.2216145,3.4800180) (3.2389818,3.5136729) (3.2524336,3.5475009) (3.2619467,3.5814156)};\draw[axis](-.12,-.19)--(2.55,4.0375)node[above]{$z$};\node[below]at(1,0){$t=0$};\node[right]at(3.3,3.7){$t=6\pi$};
\end{tikzpicture}
\caption{螺旋上方线密度更大，不能只按几何圈数判断质心。}\label{fig:c4-e2-q}
\end{figure}
\begin{solution}
图~\ref{fig:c4-e2-q}中的螺旋线转过三整圈，但线密度随高度增加。令$S=6\pi$。先求切向量
\[\vct r'(t)=\frac L{\sqrt2}(-\sin t,\cos t,1),\qquad
|\vct r'(t)|=\frac L{\sqrt2}\sqrt{\sin^2t+\cos^2t+1}=L.\]
因此
\[
 z=\frac{Lt}{\sqrt2},\qquad
 \dd m=\lambda_0L(1+t/\sqrt2)\dd t.
\]
记无量纲质量积分$D$，由
\[
 D=\int_0^S(1+t/\sqrt2)\dd t
 =\left[t+\frac{t^2}{2\sqrt2}\right]_0^S
 =S+\frac{S^2}{2\sqrt2}
\]
得$\boxed{M=\lambda_0LD}$。
\begin{figure}[H]\centering
\begin{tikzpicture}[>=Stealth,line cap=round,line join=round]
\draw[axis](0,0)--(6.7,0)node[right]{$t$};\draw[axis](0,-1.6)--(0,1.7)node[above]{$y/L$};\draw[curve,domain=0:6.283185,samples=121]plot(\x,{sin(\x r)/sqrt(2)});\draw[guide](3.141593,-1.3)--(3.141593,1.3);\node[above]at(1.5,.85){前半周较轻};\node[below]at(4.7,-.85){后半周较重};
\end{tikzpicture}
\caption{画出一圈的正负$y$分布；密度随$t$增加使后半周占更大权重。}\label{fig:c4-e2-s}
\end{figure}
图~\ref{fig:c4-e2-s}按一圈画出正负$y$区间。加权项用分部积分：
\[\int t\cos t\dd t=t\sin t-\int\sin t\dd t=t\sin t+\cos t,\]
\[\int t\sin t\dd t=-t\cos t+\int\cos t\dd t=-t\cos t+\sin t.\]
原函数中的任意常数在定积分中相消。横向一阶矩所需的四个定积分为
\[
 \begin{aligned}
 \int_0^S\cos t\dd t&=[\sin t]_0^S=0,&
 \int_0^S\sin t\dd t&=[-\cos t]_0^S=0,\\
 \int_0^St\cos t\dd t&=[t\sin t+\cos t]_0^S=0,&
 \int_0^St\sin t\dd t&=[-t\cos t+\sin t]_0^S=-S.
 \end{aligned}
\]
因此
\[
 \int x\dd m=0,\qquad
 \int y\dd m=-\frac{\lambda_0L^2S}{2},
 \qquad
 \boxed{x_C=0,\quad y_C=-\frac{LS}{2D}}.
\]
每圈的负贡献相同：图中后半周处于更高位置，密度较大，又对应$y<0$，所以$y$方向质量矩未能抵消；$x$方向的加权积分则为零。

高度方向的一阶矩为
\[
 \begin{aligned}
 \int z\dd m
 &=\frac{\lambda_0L^2}{\sqrt2}
 \int_0^S(t+t^2/\sqrt2)\dd t\\
 &=\frac{\lambda_0L^2}{\sqrt2}
 \left[\frac{t^2}{2}+\frac{t^3}{3\sqrt2}\right]_0^S\\
 &=\lambda_0L^2\left(\frac{S^2}{2\sqrt2}+\frac{S^3}{6}\right),
 \end{aligned}
\]
故
\[
 \boxed{z_C=L\frac{S^2/(2\sqrt2)+S^3/6}{D}}.
\]
若密度恒定，质心在轴线上，高度为$LS/(2\sqrt2)$。本题与该高度之差，通分后为
\[
 z_C-\frac{LS}{2\sqrt2}
 =\frac LD\left[\frac{S^2}{2\sqrt2}+\frac{S^3}{6}
 -\frac S{2\sqrt2}\left(S+\frac{S^2}{2\sqrt2}\right)\right]
 =\frac{LS^3}{24D}>0.
\]
密度随高度增加，确使质心高于均匀线的质心。
\end{solution}

\Needspace{8\baselineskip}
\begin{example}[螺旋路径上的功]\label{ex:c4-e3}
曲线$\vct r=L(\cos t,\sin t,t)$，$0\le t\le3\pi$。力场$\vct F=F_0(-x/(2L),-y/(2L),1/4)$。求做功，并判断能否改走一条更简单的路径。 见图~\ref{fig:c4-e3-q}。
\end{example}

\begin{figure}[H]\centering
\begin{tikzpicture}[>=Stealth,line cap=round,line join=round]
\draw[curve]plot coordinates{(1.0000000,0.0000000) (1.0032765,0.0169628) (1.0055666,0.0339147) (1.0068717,0.0508450) (1.0071943,0.0677427) (1.0065379,0.0845972) (1.0049067,0.1013976) (1.0023061,0.1181333) (0.9987425,0.1347937) (0.9942229,0.1513681) (0.9887556,0.1678462) (0.9823498,0.1842176) (0.9750154,0.2004719) (0.9667635,0.2165991) (0.9576058,0.2325891) (0.9475552,0.2484321) (0.9366253,0.2641182) (0.9248305,0.2796379) (0.9121863,0.2949818) (0.8987089,0.3101407) (0.8844152,0.3251054) (0.8693231,0.3398670) (0.8534513,0.3544170) (0.8368190,0.3687468) (0.8194465,0.3828481) (0.8013546,0.3967130) (0.7825648,0.4103337) (0.7630995,0.4237026) (0.7429815,0.4368124) (0.7222345,0.4496560) (0.7008826,0.4622267) (0.6789506,0.4745180) (0.6564640,0.4865236) (0.6334485,0.4982376) (0.6099307,0.5096542) (0.5859374,0.5207682) (0.5614961,0.5315744) (0.5366346,0.5420681) (0.5113812,0.5522448) (0.4857645,0.5621003) (0.4598134,0.5716308) (0.4335575,0.5808329) (0.4070262,0.5897032) (0.3802494,0.5982390) (0.3532574,0.6064377) (0.3260805,0.6142971) (0.2987491,0.6218153) (0.2712941,0.6289909) (0.2437463,0.6358226) (0.2161364,0.6423096) (0.1884956,0.6484513) (0.1608547,0.6542476) (0.1332449,0.6596987) (0.1056970,0.6648051) (0.0782420,0.6695676) (0.0509107,0.6739874) (0.0237337,0.6780660) (-0.0032583,0.6818054) (-0.0300350,0.6852076) (-0.0565663,0.6882753) (-0.0828223,0.6910113) (-0.1087733,0.6934189) (-0.1343901,0.6955014) (-0.1596435,0.6972628) (-0.1845050,0.6987071) (-0.2089463,0.6998390) (-0.2329395,0.7006631) (-0.2564574,0.7011845) (-0.2794728,0.7014085) (-0.3019595,0.7013410) (-0.3238915,0.7009878) (-0.3452434,0.7003551) (-0.3659904,0.6994495) (-0.3861083,0.6982778) (-0.4055737,0.6968469) (-0.4243634,0.6951643) (-0.4424554,0.6932375) (-0.4598279,0.6910742) (-0.4764602,0.6886824) (-0.4923320,0.6860705) (-0.5074241,0.6832469) (-0.5217178,0.6802203) (-0.5351952,0.6769995) (-0.5478394,0.6735937) (-0.5596341,0.6700120) (-0.5705641,0.6662639) (-0.5806147,0.6623590) (-0.5897724,0.6583070) (-0.5980243,0.6541179) (-0.6053587,0.6498016) (-0.6117645,0.6453683) (-0.6172318,0.6408283) (-0.6217513,0.6361919) (-0.6253150,0.6314696) (-0.6279156,0.6266719) (-0.6295468,0.6218095) (-0.6302032,0.6168931) (-0.6298806,0.6119334) (-0.6285754,0.6069412) (-0.6262854,0.6019273) (-0.6230089,0.5969026) (-0.6187455,0.5918779) (-0.6134958,0.5868640) (-0.6072611,0.5818718) (-0.6000439,0.5769121) (-0.5918477,0.5719957) (-0.5826767,0.5671333) (-0.5725363,0.5623357) (-0.5614328,0.5576134) (-0.5493734,0.5529770) (-0.5363663,0.5484369) (-0.5224206,0.5440036) (-0.5075464,0.5396873) (-0.4917547,0.5354982) (-0.4750572,0.5314462) (-0.4574667,0.5275413) (-0.4389970,0.5237932) (-0.4196624,0.5202116) (-0.3994784,0.5168057) (-0.3784611,0.5135849) (-0.3566277,0.5105583) (-0.3339958,0.5077347) (-0.3105841,0.5051228) (-0.2864120,0.5027311) (-0.2614996,0.5005677) (-0.2358679,0.4986409) (-0.2095383,0.4969583) (-0.1825331,0.4955274) (-0.1548754,0.4943557) (-0.1265885,0.4934501) (-0.0976968,0.4928174) (-0.0682250,0.4924642) (-0.0381985,0.4923967) (-0.0076432,0.4926208) (0.0234144,0.4931422) (0.0549475,0.4939662) (0.0869286,0.4950981) (0.1193299,0.4965424) (0.1521232,0.4983038) (0.1852797,0.5003864) (0.2187706,0.5027939) (0.2525664,0.5055299) (0.2866375,0.5085976) (0.3209541,0.5119999) (0.3554859,0.5157392) (0.3902027,0.5198179) (0.4250738,0.5242377) (0.4600686,0.5290001) (0.4951563,0.5341065) (0.5303060,0.5395576) (0.5654867,0.5453539) (0.6006673,0.5514956) (0.6358170,0.5579826) (0.6709047,0.5648143) (0.7058996,0.5719899) (0.7407707,0.5795081) (0.7754875,0.5873675) (0.8100193,0.5955662) (0.8443359,0.6041020) (0.8784070,0.6129724) (0.9122028,0.6221744) (0.9456936,0.6317049) (0.9788502,0.6415604) (1.0116434,0.6517371) (1.0440447,0.6622308) (1.0760258,0.6730370) (1.1075589,0.6841510) (1.1386166,0.6955676) (1.1691719,0.7072816) (1.1991984,0.7192872) (1.2286702,0.7315785) (1.2575619,0.7441492) (1.2858487,0.7569928) (1.3135065,0.7701026) (1.3405117,0.7834715) (1.3668412,0.7970922) (1.3924730,0.8109571) (1.4173853,0.8250585) (1.4415574,0.8393882) (1.4649691,0.8539382) (1.4876010,0.8686998) (1.5094345,0.8836645) (1.5304518,0.8988234) (1.5506358,0.9141673) (1.5699703,0.9296870) (1.5884401,0.9453731) (1.6060305,0.9612161) (1.6227280,0.9772061) (1.6385198,0.9933333) (1.6533940,1.0095876) (1.6673396,1.0259590) (1.6803467,1.0424371) (1.6924061,1.0590115) (1.7035096,1.0756719) (1.7136500,1.0924076) (1.7228210,1.1092080) (1.7310173,1.1260625) (1.7382345,1.1429602) (1.7444691,1.1598905) (1.7497189,1.1768424) (1.7539822,1.1938052) (1.7572587,1.2107680) (1.7595488,1.2277199) (1.7608539,1.2446502) (1.7611766,1.2615479) (1.7605201,1.2784024) (1.7588890,1.2952028) (1.7562884,1.3119385) (1.7527247,1.3285989) (1.7482051,1.3451733) (1.7427379,1.3616514) (1.7363320,1.3780228) (1.7289977,1.3942771) (1.7207457,1.4104043) (1.7115880,1.4263943) (1.7015374,1.4422373) (1.6906075,1.4579234) (1.6788128,1.4734431) (1.6661686,1.4887871) (1.6526911,1.5039459) (1.6383975,1.5189106) (1.6233054,1.5336722) (1.6074335,1.5482222) (1.5908013,1.5625520) (1.5734287,1.5766533) (1.5553368,1.5905182) (1.5365470,1.6041389) (1.5170817,1.6175078) (1.4969637,1.6306176) (1.4762167,1.6434612) (1.4548648,1.6560319) (1.4329329,1.6683232) (1.4104462,1.6803288) (1.3874307,1.6920428) (1.3639129,1.7034594) (1.3399196,1.7145734) (1.3154783,1.7253796) (1.2906168,1.7358733) (1.2653634,1.7460500) (1.2397467,1.7559055) (1.2137957,1.7654360) (1.1875397,1.7746381) (1.1610084,1.7835084) (1.1342317,1.7920442) (1.1072396,1.8002429) (1.0800627,1.8081023) (1.0527314,1.8156206) (1.0252764,1.8227961) (0.9977285,1.8296278) (0.9701186,1.8361148) (0.9424778,1.8422565) (0.9148369,1.8480528) (0.8872271,1.8535039) (0.8596792,1.8586103) (0.8322242,1.8633728) (0.8048929,1.8677926) (0.7777159,1.8718712) (0.7507239,1.8756106) (0.7239472,1.8790128) (0.6974159,1.8820805) (0.6711599,1.8848166) (0.6452089,1.8872241) (0.6195922,1.8893066) (0.5943388,1.8910680) (0.5694773,1.8925123) (0.5450360,1.8936442) (0.5210427,1.8944683) (0.4975249,1.8949897) (0.4745094,1.8952138) (0.4520227,1.8951462) (0.4300908,1.8947930) (0.4087389,1.8941603) (0.3879919,1.8932547) (0.3678739,1.8920830) (0.3484086,1.8906522) (0.3296188,1.8889695) (0.3115269,1.8870427) (0.2941543,1.8848794) (0.2775221,1.8824876) (0.2616502,1.8798757) (0.2465581,1.8770521) (0.2322645,1.8740255) (0.2187870,1.8708047) (0.2061428,1.8673989) (0.1943481,1.8638172) (0.1834182,1.8600691) (0.1733675,1.8561642) (0.1642099,1.8521122) (0.1559579,1.8479231) (0.1486236,1.8436068) (0.1422177,1.8391735) (0.1367505,1.8346335) (0.1322309,1.8299971) (0.1286672,1.8252748) (0.1260666,1.8204771) (0.1244355,1.8156147) (0.1237790,1.8106983) (0.1241017,1.8057386) (0.1254068,1.8007464) (0.1276969,1.7957326) (0.1309734,1.7907078)};\draw[axis](-.12,-.19)--(1.5,2.375)node[above]{$z$};\node[below]at(1,0){$t=0$};\node[right]at(.2,1.9){$t=3\pi$};
\end{tikzpicture}
\caption{路径绕轴上升，但场的水平功与竖直功可以分开。}\label{fig:c4-e3-q}
\end{figure}
\begin{solution}
图~\ref{fig:c4-e3-q}中的螺旋路径一边绕轴一边上升。将力与位移逐分量相乘，可以分别计算横向功和竖直功；再利用势能检验改走折线的结果。


沿曲线$\dd\vct r=L(-\sin t,\cos t,1)\dd t$，所以
\[\vct F\cdot\dd\vct r=F_0L[\tfrac12\cos t\sin t-\tfrac12\sin t\cos t+\tfrac14]\dd t.
\]
水平两项逐点抵消，故
\[W=\frac{F_0L}{4}\int_0^{3\pi}\dd t
=\frac{F_0L}{4}[t]_0^{3\pi}=\boxed{3\pi F_0L/4}.\]
势能可取$U=F_0(x^2+y^2)/(4L)-F_0z/4$。逐分量求导有
\[-\nabla U=F_0\left(-\frac{2x}{4L},-\frac{2y}{4L},\frac14\right)=\vct F.\]
因此可按图~\ref{fig:c4-e3-s}改走折线：先在底面从$(L,0,0)$到$(-L,0,0)$，再竖直上升至$(-L,0,3\pi L)$。两段功分别为
\[W_1=-\frac{F_0}{2L}\int_L^{-L}x\dd x
=-\frac{F_0}{4L}[x^2]_L^{-L}=0,\]
\[W_2=\frac{F_0}{4}\int_0^{3\pi L}\dd z
=\frac{F_0}{4}[z]_0^{3\pi L}=\frac{3\pi F_0L}{4}.\]

原路径与替代路径有相同端点，且势在全空间单值光滑，所以改路成立。单凭图形绕了几圈不能判断功的大小；力场是否保守才是改路的依据。


横向两项合为$-F_0(x\dd x+y\dd y)/(2L)$，而螺旋线上$x^2+y^2=L^2$，有$x\dd x+y\dd y=0$。所以横向力始终垂直于水平切向位移，只剩$F_0\dd z/4$做功。

相同结果也可由$W=U(A)-U(B)$求出。势能中的横向平方项在两端相等，只有高度项改变，与逐段积分的抵消关系对应。

\begin{figure}[H]\centering
\begin{tikzpicture}[>=Stealth,line cap=round,line join=round]
\draw[axis](-2,0)--(2,0)node[right]{$x$};\draw[axis](0,0)--(0,3)node[above]{$z$};\draw[axis](1.4,0)--(-1.4,0);\draw[axis](-1.4,0)--(-1.4,2.5);\fill(1.4,0)circle(1.5pt)node[below]{$A$};\fill(-1.4,0)circle(1.5pt)node[below]{$C$};\fill(-1.4,2.5)circle(1.5pt)node[above]{$B$};\node[right]at(.3,1.5){$W_{AC}=0,\ W_{CB}=3\pi F_0L/4$};
\end{tikzpicture}
\caption{保守场允许改走折线，端点高度差决定本题净功。}\label{fig:c4-e3-s}
\end{figure}
\end{solution}
\subsection{曲面面积微元}
参数曲面$\vct r(u,v)$上，两条参数线微位移张成小平行四边形，所以
\[\dd S=|\vct r_u\times\vct r_v|\dd u\dd v,\qquad
\dd\vct S=\pm(\vct r_u\times\vct r_v)\dd u\dd v.
\]
标量面积分用$\dd S$，通量用有向面积元；后者正负号由所选法向决定。对图形$z=z(x,y)$，
\[\dd S=\sqrt{1+z_x^2+z_y^2}\dd x\dd y,\quad
\dd\vct S_{\rm up}=(-z_x,-z_y,1)\dd x\dd y.
\]
若改用$z=z(r,\varphi)$，参数切向量为$\vct r_r=\vct e_r+z_r\vct e_z$、$\vct r_\varphi=r\vct e_\varphi+z_\varphi\vct e_z$，二者叉积为
\[\vct r_r\times\vct r_\varphi
=-rz_r\vct e_r-z_\varphi\vct e_\varphi+r\vct e_z.\]
由当地正交单位基矢求模，得到
\[\dd S=r\sqrt{1+z_r^2+(z_\varphi/r)^2}\dd r\dd\varphi.
\]
若曲面在球坐标中由径向图形$R=R(\theta,\varphi)>0$表示，两个参数切向量分别为$R_\theta\vct e_r+R\vct e_\theta$和$R_\varphi\vct e_r+R\sin\theta\vct e_\varphi$。逐项叉乘得
\[\vct r_\theta\times\vct r_\varphi
=R^2\sin\theta\vct e_r-RR_\theta\sin\theta\vct e_\theta-RR_\varphi\vct e_\varphi.\]
取模并提出$R^2\sin\theta$，得到
\[\boxed{\dd S=R^2\sin\theta\sqrt{1+(R_\theta/R)^2+[R_\varphi/(R\sin\theta)]^2}\dd\theta\dd\varphi}.
\]
各导数必须除以相应尺度因子，根号内才量纲一致。球面$R$为常数时恢复$R^2\sin\theta\dd\theta\dd\varphi$。这类径向参数化要求每个方向与曲面相交一次；若一个方向有多个交点，需要分片积分。

\subsection{曲面面积、投影与通量}
曲面参数线的两条小边分别为$\vct r_u\dd u$与$\vct r_v\dd v$。它们不必正交，实际面积为叉积模，等于两边长乘夹角正弦。用内积可写成
\[|\vct r_u\times\vct r_v|^2
=(\vct r_u\cdot\vct r_u)(\vct r_v\cdot\vct r_v)
-(\vct r_u\cdot\vct r_v)^2.\]
参数线正交时，面积就是两边长之积；非正交时，还须乘夹角的正弦。

对图形$z=z(x,y)$，取$\vct r_x=(1,0,z_x)$、$\vct r_y=(0,1,z_y)$，叉积为$(-z_x,-z_y,1)$。向上单位法向须除以$\sqrt{1+z_x^2+z_y^2}$，而实际面积元又乘回同一根号。因此通量中常直接使用向量面积元，不必先把法向归一化再把面积放大一次。

正的标量面积、单位法向以及投影面积分别承担不同作用。若求曲面质量，积面密度乘$\dd S$；若求穿过曲面的流量，积场的法向分量乘$\dd S$；若用平面投影计算，法向与面积的组合已吸收到$\dd\vct S$。习题~\ref{ex:c4-h4}中特别需要根据所围体积确定外法向，向上法向未必等于外法向。

球径向图形的面积元同样由两个参数矢量求叉积。在极点附近，方位角基矢退化，公式中的$1/\sin\theta$应与整体$\sin\theta$和光滑曲面极限一起理解。若曲面不是关于原点的单值径向图形，便分片参数化，再将各片面积或通量相加；单个方向有多个交点时直接套一张径向图会漏计曲面。


\section{多重积分与区域划分}\label{sec:c4-multiple}
连续分布的总量由密度与积分区域共同决定。几何条件确定区域及投影，再据此安排积分顺序；空穴、重叠和多支边界分别分区处理。

图~\ref{fig:c4-oriented-area}将两条切向微边组成的平行四边形与其法向放在一起。面积取叉积的模；通量则保留叉积的方向。交换参数顺序时，两条边所张面积不变，但法向反转。
\begin{figure}[H]\centering
\begin{tikzpicture}[>=Stealth,line cap=round,line join=round]
\coordinate(O)at(0,0);\coordinate(U)at(3,.35);\coordinate(V)at(.85,1.7);\coordinate(W)at(3.85,2.05);
\fill[gray!16](O)--(U)--(W)--(V)--cycle;\draw[curve](O)--(U)--(W)--(V)--cycle;
\draw[->,thick](O)--(U)node[midway,below]{$\vct R_u\dd u$};\draw[->,thick](O)--(V)node[midway,left]{$\vct R_v\dd v$};
\draw[->,thick](1.925,1.025)--(1.55,2.9)node[above]{$\vct n$};\node at(2.8,1.1){$\dd S$};
\node[below]at(1.9,-.65){$\vct n\dd S=(\vct R_u\times\vct R_v)\dd u\dd v$};
\end{tikzpicture}
\caption{参数曲面的面积与定向来自同一叉积，标量面积舍去方向，通量保留方向。}\label{fig:c4-oriented-area}
\end{figure}

\subsection{积分区域与积分次序}
连续函数在规则区域的重积分可以化为累次积分。例如$D=\{a\le x\le b,\ y_-(x)\le y\le y_+(x)\}$时
\[\iint_D f\dd x\dd y=\int_a^b\!\dd x\int_{y_-(x)}^{y_+(x)}f(x,y)\dd y.
\]
换次序改变的是扫描区域的方法，不改变区域本身。若水平截线穿过的边界发生变化，就必须分段，不能只把两个积分号和上下限机械对调。积分区域、密度和积分微元还需同时考虑，奇偶对称性只有在三者共同满足相应变换时才有效。

柱坐标的体积元为$r\dd r\dd\varphi\dd z$，球坐标为$r^2\sin\theta\dd r\dd\theta\dd\varphi$。二维极坐标的Jacobian为$r$：径向小边长$\dd r$，角向小边长$r\dd\varphi$。Jacobian给出参数微元对应的实际面积或体积。一般变换中使用绝对值$|\det\partial(x,y,z)/\partial(u,v,w)|$，并要求参数域除零测边界外不重复覆盖。

\Needspace{8\baselineskip}
\subsection{累次积分与区域分割}
确定累次积分时，先选外层变量，再固定它看内层允许的区间。对$y_-(x)\le y\le y_+(x)$，竖直扫描线的长度和被积函数共同决定内层积分。交换次序时，应重新用水平扫描线与原区域相交，而不是机械交换两个积分号。

对多个不等式给出的三维公共区域，固定某些坐标后，允许高度常是几个上界的最小值、几个下界的最大值。控制边界改变的位置就是需要分区的位置。例~\ref{ex:c4-e7}用固定$z$的正方形截面能一次同时满足两圆柱约束；投影到$xy$平面后，则须用$y=x$区分控制高度的圆柱。

若函数连续、区域有界且边界规则，Fubini定理允许把重积分写成累次积分并交换顺序。对于非负函数，可在更一般情形使用单调意义的交换；但含正负抵消的无穷或奇异积分，只有绝对收敛或其他明确条件成立时才可任意换序。角向变号密度的积分仍应先确认径向和角向积分可积，不能单凭一个对称方向的零积分掩盖另一方向的发散。

坐标变换还改变微元。二维极坐标的微小扇环有径向宽度$\dd r$与角向长度$r\dd\varphi$，面积因子是$r$；球坐标小体的三边为$\dd r,r\dd\theta,r\sin\theta\dd\varphi$，体积因子为$r^2\sin\theta$。被积函数中的$r$若来自密度或转动半径，仍须另乘，不能与Jacobian中的同名因子合并丢掉。

对称性应同时作用于区域、微元和被积函数。反射保持圆盘时，奇函数项可配对相消；若只保留半圆盘，区域不再对称，该项通常非零。对均匀球的各坐标平方矩相等，则来自整体旋转不变性，之后仍需通过径向积分确定每一项的具体大小。

\begin{example}[圆盘上的偶项与奇项]\label{ex:c4-e4}
在圆盘$x^2+y^2\le9$上计算$\iint_D(x^2+y)\dd x\dd y$，分别说明对称性和极坐标的作用。 见图~\ref{fig:c4-e4-q}。
\end{example}

\begin{figure}[H]\centering
\begin{tikzpicture}[>=Stealth,line cap=round,line join=round]
\fill[gray!12](0,0)circle(1.5);\draw[curve](0,0)circle(1.5);\draw[axis](-1.9,0)--(2,0)node[right]{$x$};\draw[axis](0,-1.8)--(0,2)node[above]{$y$};\node at(.9,.9){$R=3$};
\end{tikzpicture}
\caption{圆盘关于两坐标轴对称，$y$项对上、下半区反号。}\label{fig:c4-e4-q}
\end{figure}
\begin{solution}
图~\ref{fig:c4-e4-q}中的圆盘及面积元都在$x$轴反射下保持不变，被积函数中的$y$却反号。配对消去这一项后，余下$x^2$用极坐标积分。


反射$(x,y)\mapsto(x,-y)$保持圆盘和面积元不变，$y$项反号，所以$\iint_Dy\dd A=0$。对$x^2$项，图~\ref{fig:c4-e4-s}的极坐标给出$x=r\cos\varphi$、$\dd A=r\dd r\dd\varphi$，
\[I=\int_0^3r^3\dd r\int_0^{2\pi}\cos^2\varphi\dd\varphi
\]
中径向的$r^3$由$x^2$中的$r^2$和面积元中的$r$共同产生。两个积分分别为
\[\int_0^3r^3\dd r=\left[\frac{r^4}{4}\right]_0^3=\frac{81}4,\]
\[\int_0^{2\pi}\cos^2\varphi\dd\varphi
=\int_0^{2\pi}\frac{1+\cos2\varphi}{2}\dd\varphi
=\left[\frac\varphi2+\frac{\sin2\varphi}{4}\right]_0^{2\pi}=\pi.\]
所以$\boxed{I=81\pi/4}$。

还可用图~\ref{fig:c4-e4-s}中的旋转对称关系独立计算。记两个平方矩为$I_x=\iint_Dx^2\dd A$、$I_y=\iint_Dy^2\dd A$，则$I_x=I_y$，而
\[I_x+I_y=\iint_Dr^2\dd A
=2\pi\int_0^3r^3\dd r=\frac{81\pi}{2}.\]
每项取一半，也得到$I_x=81\pi/4$。反射对称与旋转对称的作用不同：若改成上半圆盘，$y$项不再抵消；若只保留一般扇区，两个平方矩也未必相等。

\begin{figure}[H]\centering
\begin{tikzpicture}[>=Stealth,line cap=round,line join=round]
\draw[curve](0,0)circle(1.5);\draw[axis](0,0)--(1.2,.9)node[midway,above]{$r$};\draw[curve](.7,0)arc(0:36.87:.7);\node at(.9,.25){$\varphi$};\fill(1,.5)circle(1.5pt)node[right]{$+y$};\fill(1,-.5)circle(1.5pt)node[right]{$-y$};\draw[guide](1,.5)--(1,-.5);
\end{tikzpicture}
\caption{反射配对消去$y$项，余下偶项用极坐标积分。}\label{fig:c4-e4-s}
\end{figure}
\end{solution}

\Needspace{8\baselineskip}
\begin{example}[旋转抛物面的面积]\label{ex:c4-e5}
求抛物面$z=k(x^2+y^2)$在$x^2+y^2\le r_0^2$内的曲面面积，$k>0$。写出直角坐标与极坐标形式，并讨论$k\to0$。 见图~\ref{fig:c4-e5-q}。
\end{example}

\begin{figure}[H]\centering
\begin{tikzpicture}[>=Stealth,line cap=round,line join=round]
\draw[axis](-2,0)--(2,0)node[right]{$x$};\draw[axis](0,-.2)--(0,2.5)node[above]{$z$};\draw[curve,domain=-1.6:1.6]plot(\x,{.6*(\x)^2});\draw[dashed](-1.6,1.536)--(1.6,1.536);\node[right]at(1.6,1.536){$|x|=r_0$};
\end{tikzpicture}
\caption{子午截面$y=0$中的正半支旋转生成曲面，圆环微元含斜率修正。}\label{fig:c4-e5-q}
\end{figure}
\begin{solution}
曲面斜率使实际面积大于水平投影面积。图~\ref{fig:c4-e5-q}的子午线在半径$r$处斜率为$2kr$，因此应把投影圆环的宽度$\dd r$换成曲面上的斜宽$\dd s$；周向长度$2\pi r$则保持原值。


因为$z_x=2kx,z_y=2ky$，面积为
\[S=\iint_{x^2+y^2\le r_0^2}\sqrt{1+4k^2(x^2+y^2)}\dd x\dd y.
\]
图~\ref{fig:c4-e5-s}取半径$r$处的窄圆环，其斜长$\dd s=\sqrt{1+4k^2r^2}\dd r$、周长$2\pi r$，所以
\[S=2\pi\int_0^{r_0}r\sqrt{1+4k^2r^2}\dd r.\]
令$u=1+4k^2r^2$，则$\dd u=8k^2r\dd r$，积分上下限为$1$与$1+4k^2r_0^2$。于是
\[\begin{aligned}
S&=\frac\pi{4k^2}\int_1^{1+4k^2r_0^2}u^{1/2}\dd u\\
&=\frac\pi{4k^2}\left[\frac23u^{3/2}\right]_1^{1+4k^2r_0^2}
=\boxed{\frac\pi{6k^2}[(1+4k^2r_0^2)^{3/2}-1]}.
\end{aligned}\]
$k\to0$时不能把式子中的分母直接置零。展开$(1+u)^{3/2}=1+3u/2+O(u^2)$，得$S\to\pi r_0^2$，正是平面圆盘面积。对任意$k>0$，被积根号大于$1$，曲面面积严格大于底面投影面积。


若需要近平面时的面积修正，令$z=4k^2r_0^2\ll1$。将二项式保留到$z^3$，即
\[
 (1+z)^{3/2}-1=\frac32z+\frac38z^2-\frac1{16}z^3+O(z^4).
\]
依次代入$z$的幂，并乘上$\pi/(6k^2)$，各项系数为
\[
 \frac{\pi}{6k^2}\frac32(4k^2r_0^2)=\pi r_0^2,\qquad
 \frac{\pi}{6k^2}\frac38(16k^4r_0^4)=\pi k^2r_0^4,
\]
\[
 -\frac{\pi}{6k^2}\frac1{16}(64k^6r_0^6)
 =-\frac23\pi k^4r_0^6.
\]
所以
\[S=\pi r_0^2+\pi k^2r_0^4-\frac23\pi k^4r_0^6+O(k^6r_0^8).\]
图中斜长与水平宽度的差从斜率平方阶开始，所以首项面积修正为正且与$k^2$成正比。

另一种参数化是$\vct R(r,\varphi)=(r\cos\varphi,r\sin\varphi,kr^2)$。其切向量为
\[\vct R_r=(\cos\varphi,\sin\varphi,2kr),\qquad
\vct R_\varphi=(-r\sin\varphi,r\cos\varphi,0).\]
二者点积为零，长度分别为$\sqrt{1+4k^2r^2}$与$r$，所以
$|\vct R_r\times\vct R_\varphi|=r\sqrt{1+4k^2r^2}$，与圆环法给出的面积系数一致。

图~\ref{fig:c4-e5-s}中的三角形应理解为微分极限；有限宽圆环的弦长并不严格等于曲线弧长。积分把窄环宽度趋于零之后的首阶面积累加，所以不需要额外补上“曲率修正项”。量纲上$k$的单位是长度的倒数，答案中的$1/k^2$与面积单位相符。

\begin{figure}[H]\centering
\begin{tikzpicture}[>=Stealth,line cap=round,line join=round]
\draw[axis](0,0)--(3,0)node[right]{$r$};\draw[axis](0,0)--(0,2.8)node[above]{$z$};\draw[curve,domain=0:2]plot(\x,{.5*(\x)^2});\draw[guide](1.1,.605)--(1.35,.605)--(1.35,.88);\draw[thick](1.1,.605)--(1.35,.88)node[midway,above left]{$\dd s$};\node[below]at(1.225,.605){$\dd r$};\node[right]at(1.35,.74){$2kr\dd r$};
\end{tikzpicture}
\caption{圆环面积为$2\pi r\dd s$，不是投影环面积$2\pi r\dd r$。}\label{fig:c4-e5-s}
\end{figure}
\end{solution}

\Needspace{8\baselineskip}
\begin{example}[球与圆锥的体积和质心]\label{ex:c4-e6}
用柱坐标与球坐标求半径$R$的球体体积，并求球面面积。另有均匀实心圆锥$r\le kz,0\le z\le h$，求其质量与质心高度。 见图~\ref{fig:c4-e6-q}。
\end{example}

\begin{figure}[H]\centering
\begin{tikzpicture}[>=Stealth,line cap=round,line join=round]
\begin{scope}\draw[curve](0,0)circle(1.2);\draw[dashed](0,.5)ellipse(1.09 and .25);\draw[axis](0,-1.5)--(0,1.6)node[above]{$z$};\node at(0,-1.9){球};\end{scope}\begin{scope}[shift={(4,0)}]\draw[curve](0,-1)--(-1.2,1.1)--(1.2,1.1)--cycle;\draw[dashed](0,0)ellipse(.571 and .16);\draw[axis](0,-1.3)--(0,1.6)node[above]{$z$};\node at(0,-1.9){圆锥};\end{scope}
\end{tikzpicture}
\caption{水平薄片的截面积决定体积与质量权重。}\label{fig:c4-e6-q}
\end{figure}
\begin{solution}
球和圆锥都具有轴对称性，但合适的分层变量不同。图~\ref{fig:c4-e6-q}中球可以选同轴圆柱壳或球壳；圆锥选水平圆盘后，半径与高度成正比，体积和一阶矩都化成幂函数积分。


球体在轴距$r$处的高度为$2\sqrt{R^2-r^2}$，柱坐标积分为
\[V=\int_0^{2\pi}\dd\varphi\int_0^R2r\sqrt{R^2-r^2}\dd r
=4\pi\int_0^Rr\sqrt{R^2-r^2}\dd r.\]
令$u=R^2-r^2$，则$\dd u=-2r\dd r$，$r=0,R$对应$u=R^2,0$，因此
\[V=-2\pi\int_{R^2}^0u^{1/2}\dd u
=2\pi\left[\frac23u^{3/2}\right]_0^{R^2}
=\frac{4\pi R^3}{3}.\]
球坐标则给出
\[\begin{aligned}
V&=\int_0^Rr^2\dd r\int_0^\pi\sin\theta\dd\theta
\int_0^{2\pi}\dd\varphi\\
&=\left[\frac{r^3}{3}\right]_0^R[-\cos\theta]_0^\pi
[\varphi]_0^{2\pi}=\frac{R^3}{3}\cdot2\cdot2\pi.
\end{aligned}\]
球面固定$r=R$，面积元为$R^2\sin\theta\dd\theta\dd\varphi$，所以
\[S=R^2[-\cos\theta]_0^\pi[\varphi]_0^{2\pi}=4\pi R^2.\]

圆锥在高度$z$的截面半径$kz$，图~\ref{fig:c4-e6-s}的薄片质量$\dd m=\rho\pi k^2z^2\dd z$，因此
\[M=\rho\pi k^2\int_0^hz^2\dd z
=\rho\pi k^2\left[\frac{z^3}{3}\right]_0^h
=\frac{\rho\pi k^2h^3}{3},\]
\[\int z\dd m=\rho\pi k^2\int_0^hz^3\dd z
=\rho\pi k^2\left[\frac{z^4}{4}\right]_0^h
=\frac{\rho\pi k^2h^4}{4}.\]
取一阶矩与质量之比，$z_C=(h^4/4)/(h^3/3)=\boxed{3h/4}$。
这里高度从锥尖量起，故距底面为$h/4$。三重积分形式为
\[M=\rho\int_0^h\dd z\int_0^{2\pi}\dd\varphi\int_0^{kz}r\dd r.\]
一阶矩在被积函数中多乘一个$z$。层积分与三重积分是同一体积元经不同次序合并后的写法。


球面面积也可由薄球壳关系$\dd V=S\dd R$得到：
\[S=\frac{\dd}{\dd R}\left(\frac{4\pi R^3}{3}\right)
=\frac{4\pi}{3}\cdot3R^2=4\pi R^2.\]
图~\ref{fig:c4-e6-s}把圆锥的原点放在锥尖，因此$z_C=3h/4$与距底面$h/4$对应同一点。

\begin{figure}[H]\centering
\begin{tikzpicture}[>=Stealth,line cap=round,line join=round]
\draw[curve](0,0)--(-1.5,2.5)--(1.5,2.5)--cycle;\draw[axis](0,-.2)--(0,3)node[above]{$z$};\draw[curve](-.9,1.5)--(.9,1.5);\node[right]at(.9,1.5){$\dd m=\rho\pi k^2z^2\dd z$};\fill(0,1.875)circle(2pt)node[left]{$C:3h/4$};
\end{tikzpicture}
\caption{截面积随$z^2$增加，质心比半高位置更靠近底面。}\label{fig:c4-e6-s}
\end{figure}
\end{solution}

\Needspace{8\baselineskip}
\begin{example}[两个垂直圆柱的公共体积]\label{ex:c4-e7}
两个无限长圆柱$x^2+z^2\le r^2$、$y^2+z^2\le r^2$的轴线正交且相交。求公共部分体积，并给出投影分区的积分形式。 见图~\ref{fig:c4-e7-q}。
\end{example}

\begin{figure}[H]\centering
\begin{tikzpicture}[>=Stealth,line cap=round,line join=round]
\fill[gray!15](-1,-1)rectangle(1,1);\draw[curve](-2,-1)--(2,-1) (-2,1)--(2,1);\draw[curve](-1,-2)--(-1,2) (1,-2)--(1,2);\draw[axis](-2.3,0)--(2.5,0)node[right]{$x$};\draw[axis](0,-2.3)--(0,2.5)node[above]{$y$};\node[above]at(1,1){$(r,r)$};\node at(3,1.5){公共投影};
\end{tikzpicture}
\caption{两个圆柱在$xy$平面的投影为正交条带，公共投影是边长$2r$的正方形。}\label{fig:c4-e7-q}
\end{figure}
\begin{solution}
图~\ref{fig:c4-e7-q}中的正方形是公共体积的投影，其上的高度同时受两圆柱限制。先固定$z$，两约束便分别控制$x$和$y$，截面直接化为正方形。


固定$z$后，两个不等式给出$|x|,|y|\le\sqrt{r^2-z^2}$，所以水平截面是边长$2\sqrt{r^2-z^2}$的正方形。图~\ref{fig:c4-e7-s}中的截面积为$4(r^2-z^2)$，
\[V=\int_{-r}^r4(r^2-z^2)\dd z
=4\left[r^2z-\frac{z^3}{3}\right]_{-r}^r
=8\left(r^3-\frac{r^3}{3}\right)=\boxed{16r^3/3}.\]
若先向$xy$平面投影，须取
\[|z|\le\min\{\sqrt{r^2-x^2},\sqrt{r^2-y^2}\}.
\]
在第一象限，以$y=x$分区：$0\le y\le x\le r$时第一项较小；另一半交换$x,y$。计入上下与四象限，
\[V=16\int_0^r\dd x\int_0^x\sqrt{r^2-x^2}\dd y
=16\int_0^rx\sqrt{r^2-x^2}\dd x.
\]
内层积分是$[y\sqrt{r^2-x^2}]_0^x=x\sqrt{r^2-x^2}$。外层令$u=r^2-x^2$，则$\dd u=-2x\dd x$，两端$x=0,r$对应$u=r^2,0$，故
\[V=-8\int_{r^2}^0u^{1/2}\dd u
=8\left[\frac23u^{3/2}\right]_0^{r^2}=\frac{16r^3}{3}.\]
不能在整个正方形里只选某一个圆柱给出的高度，否则会把另一圆柱以外的部分计入。相反，固定$z$扫描天然同时满足两个约束，计算更短。


投影法前的系数$16$依次来自四个象限、第一象限中关于对角线对称的两块区域，以及上下两个半体。每个因子都对应一个几何对称，不能再把水平截面中的四倍面积重复乘入。

\begin{figure}[H]\centering
\begin{tikzpicture}[>=Stealth,line cap=round,line join=round]
\begin{scope}\draw[curve](-1,-1)rectangle(1,1);\draw[<->](-1,-1.3)--(1,-1.3)node[midway,below]{$2\sqrt{r^2-z^2}$};\node at(0,1.5){固定$z$的截面};\end{scope}\begin{scope}[shift={(4,0)}]\fill[gray!15](0,0)--(1.8,0)--(1.8,1.8)--cycle;\draw[curve](0,0)rectangle(1.8,1.8);\draw[dashed](0,0)--(1.8,1.8);\node[right]at(1.8,.7){$y<x$};\node at(.6,1.5){$y>x$};\end{scope}
\end{tikzpicture}
\caption{左图用正方形截片；右图用投影区域中不同的高度上界。}\label{fig:c4-e7-s}
\end{figure}
\end{solution}

\Needspace{8\baselineskip}
\begin{example}[角向变号电荷密度的偶极矩]\label{ex:c4-e8}
球$r\le r_0$内$\rho=\rho_0(r/r_0)\cos\theta\cos\varphi$，$\theta$为极角，$\varphi$为方位角。求总电荷及电偶极矩。 见图~\ref{fig:c4-e8-q}。
\end{example}

\begin{figure}[H]\centering
\begin{tikzpicture}[>=Stealth,line cap=round,line join=round]
\draw[curve](0,0)circle(1.6);\draw[axis](0,-1.95)--(0,2.1)node[above]{$z$};
\draw[axis](-1.95,0)--(2.15,0)node[right]{$x$};
\draw[axis](0,0)--(1.32,.9)node[above right]{$\vct r$};
\draw(0,.48)arc(90:34.3:.48);\node at(.27,.57){$\theta$};
\node at(-.85,.75){$\rho<0$};\draw[thin](.55,.95)--(1.4,1.7)--(2,1.7);\node[right]at(2,1.7){$\rho>0$};
\node[below]at(0,-2.25){$y=0$截面，示意$\rho_0>0$};
\end{tikzpicture}
\caption{极角从正$z$轴量起；$xz$截面中右侧取$\varphi=0$，左侧取$\varphi=\pi$，两侧密度异号。}\label{fig:c4-e8-q}
\end{figure}
\begin{solution}
总电荷积分密度，偶极矩则分别多乘$x,y,z$。图~\ref{fig:c4-e8-q}给出密度的正负区域；对这四个被积函数分别作反射，确定各自的抵消关系。


体积元为$r^2\sin\theta\dd r\dd\theta\dd\varphi$。总电荷中的方位角积分为
\[\int_0^{2\pi}\cos\varphi\dd\varphi=[\sin\varphi]_0^{2\pi}=0,\]
所以$Q=0$。

图~\ref{fig:c4-e8-s}按三个坐标投影求偶极矩$\vct p=\int\vct r\rho\dd V$。因为$x=r\sin\theta\cos\varphi$，
\[p_x=\frac{\rho_0}{r_0}\int_0^{r_0}r^4\dd r
\int_0^\pi\sin^2\theta\cos\theta\dd\theta
\int_0^{2\pi}\cos^2\varphi\dd\varphi=0.
\]
这里极角积分在$\theta\mapsto\pi-\theta$下反号，也可直接由原函数计算：
\[\int_0^\pi\sin^2\theta\cos\theta\dd\theta
=\left[\frac{\sin^3\theta}{3}\right]_0^\pi=0.\]
$p_y$中方位角积分为
\[\int_0^{2\pi}\sin\varphi\cos\varphi\dd\varphi
=\left[\frac{\sin^2\varphi}{2}\right]_0^{2\pi}=0;\]
$p_z$中的方位角积分则与总电荷相同。因此
\[\boxed{Q=0,\qquad\vct p=\vct0}.\]
电荷分布仍然存在，只是总电荷及三个一阶矩恰好抵消。极点处的$\varphi$不唯一，可任意补定义这些零测点的密度值，积分结果不受影响。


若对$x$分量令$u=\sin\theta$换元，必须在$\pi/2$处分段，因为$u$在整个极角区间不单调；直接识别原函数则没有这个分段问题。

图~\ref{fig:c4-e8-s}中沿$z$反射会把密度反号，却保持$x$不变，故$x\rho$成对抵消；沿$y$反射保持密度而改变$y$，故$y\rho$抵消；沿$x$反射改变密度而保持$z$，故$z\rho$抵消。这给出了三个分量各自独立的几何理由。总电荷和偶极矩为零不意味着远场处处为零，更高阶矩仍可能保留。

\begin{figure}[H]\centering
\begin{tikzpicture}[>=Stealth,line cap=round,line join=round]
\draw[curve](0,0)circle(1.5);\draw[axis](-1.9,0)--(2,0)node[right]{$x$};\draw[axis](0,-1.9)--(0,2)node[above]{$z$};\node at(.7,.7){$+$};\node at(-.7,.7){$-$};\node at(.7,-.7){$-$};\node at(-.7,-.7){$+$};\draw[guide](.9,.8)--(.9,-.8);
\end{tikzpicture}
\caption{子午截面中的四区符号显示：净电荷与偶极矩可同时为零。}\label{fig:c4-e8-s}
\end{figure}

还可算出一个未抵消的二阶矩。图~\ref{fig:c4-e8-s}四区的密度符号与$xz$的符号相同（$\rho_0>0$），所以乘积$xz\rho$在各区同号。所需径向和角向积分为
\[\int_0^{r_0}r^5\dd r=\left[\frac{r^6}{6}\right]_0^{r_0}=\frac{r_0^6}{6},\]
\[\int_0^\pi\sin^2\theta\cos^2\theta\dd\theta
=\int_0^\pi\frac{1-\cos4\theta}{8}\dd\theta
=\left[\frac\theta8-\frac{\sin4\theta}{32}\right]_0^\pi=\frac\pi8,\]
\[\int_0^{2\pi}\cos^2\varphi\dd\varphi
=\left[\frac\varphi2+\frac{\sin2\varphi}{4}\right]_0^{2\pi}=\pi.\]
代入球坐标，
\[\begin{aligned}
M_{xz}&=\int xz\rho\dd V\\
&=\frac{\rho_0}{r_0}\int_0^{r_0}r^5\dd r
\int_0^\pi\sin^2\theta\cos^2\theta\dd\theta
\int_0^{2\pi}\cos^2\varphi\dd\varphi\\
&=\frac{\rho_0}{r_0}\frac{r_0^6}{6}\frac\pi8\pi
=\boxed{\frac{\rho_0\pi^2r_0^5}{48}}\ne0.
\end{aligned}\]
这里$M_{xz}$是原始二阶矩。若采用四极矩张量定义
\[Q_{ij}=\int\rho(3x_ix_j-r^2\delta_{ij})\dd V,\]
则$Q_{xz}=Q_{zx}=3M_{xz}$；其余分量由相应反射对称性为零。对远处$\vct R=R\vct n$、$R\gg r_0$，令
$u=-2\vct n\cdot\vct r/R+r^2/R^2$，利用$(1+u)^{-1/2}=1-u/2+3u^2/8+O(u^3)$，有
\[\frac1{|\vct R-\vct r|}
=\frac1R+\frac{\vct n\cdot\vct r}{R^2}
+\frac{3(\vct n\cdot\vct r)^2-r^2}{2R^3}+O(R^{-4}).\]
乘密度后对电荷体积分，前两项由$Q=0$、$\vct p=0$消失，得到
\[\Phi(\vct R)=\frac1{4\pi\epsilon_0}
\frac{Q_{ij}n_in_j}{2R^3}+O(R^{-4})
=\frac1{4\pi\epsilon_0}\frac{3M_{xz}n_xn_z}{R^3}+O(R^{-4}).\]
因此一般方向上远势的首项按$R^{-3}$衰减。净电荷和偶极矩为零，只消去了前两阶，并未使整个外场为零。此处求和包含$xz$和$zx$两项，正好与展开式中的$1/2$抵消；系数因而由已给定的张量定义确定。

\end{solution}

\section{把边界积分化为内部积分}\label{sec:c4-theorems}
边界与内部的积分关系可以减少计算，也可能引入使用条件。选定定理前，先核实区域中的光滑性、边界方向和奇异源是否已经计入。

\subsection{Green公式与正向边界}
设$D$为平面区域，边界分段光滑，$P,Q$有连续一阶偏导。正向指沿边界行进时区域在左手侧。对外边界是逆时针，对孔洞边界则是顺时针。Green公式为
\[\boxed{\oint_{\partial D}(P\dd x+Q\dd y)=\iint_D(Q_x-P_y)\dd A}.\]
对上下边界能写成函数的区域，沿下边界向右、上边界向左，$P\dd x$的总和为$\int[P(x,y_-)-P(x,y_+)]\dd x=-\iint P_y\dd A$；对左右边界同理得到$\oint Q\dd y=\iint Q_x\dd A$。一般区域分成小块，内部公共边反向抵消，留下外边界。这个抵消过程同时解释了定理和边界定向。

有孔区域的内边界应取顺时针，使保留区域仍在左侧。将区域分成小块时，相邻公共边方向相反而抵消，孔洞边界却没有另一块区域在其内部来抵消，所以必须保留。若边界曲线自交或区域有多个连通部分，应先明确所取的各片区域及其边界，不能直接把一条复杂曲线的投影当作单一正向围域。

例~\ref{ex:c4-e9}的直接参数化与区域积分提供两种计算路径。前者逐段追踪边界方向，后者把两条边之间的厚度放入内层积分限；按相同的正向边界计算，两种方法得到同一环量。

图~\ref{fig:c4-boundary-hole}把有孔区域的边界方向明确画出。沿外边界逆时针、内边界顺时针行进，灰色积分区域都在左侧。孔内若存在场奇点，不能把孔直接补上后仍假定场在整个内部光滑。
\begin{figure}[H]\centering
\begin{tikzpicture}[>=Stealth,line cap=round,line join=round]
\fill[gray!15,even odd rule](0,0)circle(1.7)(0,0)circle(.65);\draw[curve](0,0)circle(1.7);\draw[curve](0,0)circle(.65);
\draw[->,thick](1.7,0)arc(0:65:1.7);\draw[->,thick](0,.65)arc(90:0:.65);
\draw(.6,.25)--(1.85,.75)--(2.1,.75);\node[right]at(2.1,.75){内边界：顺时针};\draw(1.5,-.8)--(2.1,-.8);\node[right]at(2.1,-.8){外边界：逆时针};\node at(-1.1,0){$D$};\node at(0,0){孔};
\end{tikzpicture}
\caption{Green公式中有孔区域的正向边界；两条边界共同组成$\partial D$。}\label{fig:c4-boundary-hole}
\end{figure}

\subsection{Gauss与Stokes公式}
对光滑向量场，散度定理写为
\[\iint_{\partial V}\vct A\cdot\vct n\dd S=\iiint_V\nabla\cdot\vct A\dd V.
\]
体积小块公共面的法向相反，内部通量相消，只留下整个体积的外表面。若区域有空腔，空腔表面也属于边界，外法向指向空腔。

Stokes公式为
\[\oint_{\partial S}\vct A\cdot\dd\vct l=\iint_S(\nabla\times\vct A)\cdot\vct n\dd S.
\]
小曲面片公共边反向抵消，最后仅留总边界。曲面法向与回路正向按右手规则配对。允许选择任何同边界的曲面，前提是场在所用曲面附近有足够正则性；不能让替代曲面穿过场的奇点。

将局部电磁方程代入两定理，得到$\iint\vct D\cdot\dd\vct S=Q_{\rm enclosed}$、$\iint\vct B\cdot\dd\vct S=0$，以及
\[\oint\vct E\cdot\dd\vct l=-\frac{\dd}{\dd t}\iint\vct B\cdot\dd\vct S,\qquad
\oint\vct H\cdot\dd\vct l=I_{\rm through}+\frac{\dd}{\dd t}\iint\vct D\cdot\dd\vct S.
\]
这里曲线与曲面在空间固定。若回路运动，磁通时间导数还包含运动效应，应结合洛伦兹力另作推导。

\Needspace{8\baselineskip}
Gauss定理的外法向是相对于所选体积确定的。实心球的外表面向外，球壳的内表面却指向被挖去的球心区域。若给定曲面的法向与闭面外法向相反，使用定理后须以负号转换回题目方向。用辅助面封口时，应同时写出所围体积的坐标范围，法向判断就有明确依据。

Stokes定理的边界方向则由曲面法向决定。用同一条有向闭路跨接不同光滑曲面时，各曲面选择与闭路配对的法向，其旋度通量相等。若把两曲面合成闭面，必须把其中一面的方向反转；此时才是两个外向通量相加为零。习题~\ref{ex:c4-h5}用圆盘和上半球展示这两种符号表述的关系。

替代曲面还须留在场的光滑定义域中。有孔涡旋场的边界圆周环量非零，若用穿过奇轴的圆盘计算其零旋度通量，就违反了Stokes定理的前提。替代曲面须与原曲线具有相同有向边界，并完全位于场的光滑定义域内。

\begin{example}[两条曲线围成区域的环量]\label{ex:c4-e9}
区域由$y=x^2$与$y=x$在$0\le x\le1$围成。边界取正向，求$\oint[y^3\dd x+(x^3+3xy^2)\dd y]$，并以直接参数化复算。 见图~\ref{fig:c4-e9-q}。
\end{example}

\begin{figure}[H]\centering
\begin{tikzpicture}[>=Stealth,line cap=round,line join=round]
\fill[gray!12](0,0)--plot[domain=0:2]({\x},{\x})--plot[domain=2:0]({\x},{(\x)^2/2})--cycle;\draw[axis](0,0)--(2.6,0)node[right]{$x$};\draw[axis](0,0)--(0,2.6)node[above]{$y$};\draw[curve,domain=0:2]plot(\x,{(\x)^2/2});\draw[curve](0,0)--(2,2);\draw[axis](.9,.405)--(1.1,.605);\draw[axis](1.1,1.1)--(.9,.9);\node[below]at(2,0){$1$};\node[left]at(0,2){$1$};
\end{tikzpicture}
\caption{下边界向右、上边界向左，区域始终在行进方向左侧。}\label{fig:c4-e9-q}
\end{figure}
\begin{solution}
图~\ref{fig:c4-e9-q}的正向围域要求先沿下边界$y=x^2$从原点到$(1,1)$，再沿直线返回。直接积分时每段都要保留这个方向；用Green公式时，方向则被统一放入$Q_x-P_y$的符号中。


令$P=y^3,Q=x^3+3xy^2$，则$Q_x-P_y=(3x^2+3y^2)-3y^2=3x^2$。图~\ref{fig:c4-e9-s}的竖条上界$x$、下界$x^2$，因此
\[\begin{aligned}
I&=\int_0^1[3x^2y]_{x^2}^{x}\dd x
=3\int_0^1(x^3-x^4)\dd x\\
&=3\left[\frac{x^4}{4}-\frac{x^5}{5}\right]_0^1
=3\left(\frac14-\frac15\right)=\boxed{3/20}.
\end{aligned}\]

直接计算时，下边界$C_1$取$x=t,y=t^2,0\le t\le1$，因此$\dd x=\dd t,\dd y=2t\dd t$，逐项代入得
\[y^3\dd x=t^6\dd t,\qquad
(x^3+3xy^2)\dd y=(t^3+3t^5)2t\dd t.\]
这段积分为
\[I_1=\int_0^1(7t^6+2t^4)\dd t
=\left[t^7+\frac{2t^5}{5}\right]_0^1=\frac75.\]
上边界$C_2$取$x=y=t$，由$(1,1)$返回$O$，故$t:1\to0$，两微分均为$\dd t$，从而
\[I_2=\int_1^0(t^3+4t^3)\dd t
=\left[\frac{5t^4}{4}\right]_1^0=-\frac54.\]
相加得$I_1+I_2=7/5-5/4=3/20$，与区域积分相同。图~\ref{fig:c4-e9-s}的两段箭头对应这两套积分限，返回段的负号已经包含在上下限中。若全部边界反向，结果为$-3/20$。

\begin{figure}[H]\centering
\begin{tikzpicture}[>=Stealth,line cap=round,line join=round]
\draw[axis](0,0)--(2.5,0)node[right]{$x$};\draw[axis](0,0)--(0,2.5)node[above]{$y$};\draw[curve](0,0)--(2,2);\draw[curve,domain=0:2]plot(\x,{(\x)^2/2});\draw[thick](1.2,.72)--(1.2,1.2);\node[right]at(1.25,1.25){$y=x$};\node[right]at(1.25,.55){$y=x^2$};\node at(3.8,1.8){$\dd A=(x-x^2)\dd x$};
\end{tikzpicture}
\caption{竖条高度必须取两条边界之差，曲线积分的结果可由两种方法互证。}\label{fig:c4-e9-s}
\end{figure}
\end{solution}

\begin{exercise}[密度向边缘减小的带电球]\label{ex:c4-v1}

半径$R$的球内体电荷密度为$\rho(r)=\rho_0(1-r^2/R^2)$，$\rho_0>0$，球外无电荷，也无表面电荷。真空介电常数为$\epsilon_0$，取$\Phi(\infty)=0$。求总电荷、全空间电场和电势，核对中心及球面边界，并确定场强最大的位置。 见图~\ref{fig:c4-v1-q}。
\end{exercise}

\begin{figure}[H]\centering
\begin{tikzpicture}[>=Stealth,line cap=round,line join=round]

\draw[curve](0,0)circle(1.4);\draw[guide](0,0)circle(.85);\fill(0,0)circle(1.5pt)node[below left]{$O$};
\draw[axis](0,0)--(.85,0)node[midway,above]{$r$};\draw[guide](0,0)--(-.99,.99)node[midway,above left]{$R$};\draw[axis](.6,.6)--(1.05,1.05)node[above right]{$\vct n$};
\begin{scope}[shift={(4.2,0)},x=2cm,y=1.6cm]
\draw[axis](0,0)--(1.3,0)node[right]{$r/R$};\draw[axis](0,0)--(0,1.3)node[above]{$\rho/\rho_0$};
\draw[curve,domain=0:1,samples=81]plot(\x,{1-(\x)^2});\node[below]at(1,0){$1$};
\end{scope}

\end{tikzpicture}
\caption{左图虚线为内部Gauss球面，右图是实际径向密度；球面处体密度连续降至零。}\label{fig:c4-v1-q}
\end{figure}
\begin{solution}
图~\ref{fig:c4-v1-q}中的密度只依赖半径，电场必沿径向。先积分内部电荷：
\[Q(r)=4\pi\int_0^r\rho(s)s^2\dd s
=4\pi\rho_0\left[\frac{s^3}{3}-\frac{s^5}{5R^2}\right]_0^r
=4\pi\rho_0\left(\frac{r^3}{3}-\frac{r^5}{5R^2}\right),\quad r\le R.
\]
总电荷为$\boxed{Q=8\pi\rho_0R^3/15}$。在图示Gauss球面上$\vct E\cdot\vct n=E_r$且大小恒定，由$4\pi r^2E_r=Q_{\rm enclosed}/\epsilon_0$，
\[\boxed{E_r(r)=\frac{\rho_0}{\epsilon_0}
\begin{cases}
r/3-r^3/(5R^2),&0\le r\le R,\\
2R^3/(15r^2),&r\ge R.
\end{cases}}\]
图~\ref{fig:c4-v1-s}上图画出这两个分支，场强由中心的零先增加再减小，不能只凭包围电荷单调增加便认定场强也单调增加。

球外势由无穷远边界确定：
\[\Phi_{\rm out}(r)=\int_r^\infty\frac{Q}{4\pi\epsilon_0s^2}\dd s
=\frac{Q}{4\pi\epsilon_0}\left[-\frac1s\right]_r^\infty
=\frac{Q}{4\pi\epsilon_0r}.\]
内部从球面向内积分：
\[\begin{aligned}
\Phi_{\rm in}(r)&=\Phi(R)+\int_r^R E_r(s)\dd s\\
&=\frac{\rho_0}{\epsilon_0}\left\{\frac{2R^2}{15}
+\left[\frac{s^2}{6}-\frac{s^4}{20R^2}\right]_r^R\right\}\\
&=\frac{\rho_0}{\epsilon_0}\left(\frac{R^2}{4}-\frac{r^2}{6}
+\frac{r^4}{20R^2}\right).
\end{aligned}\]
因此
\[\boxed{\Phi(r)=\frac{\rho_0}{\epsilon_0}
\begin{cases}
R^2/4-r^2/6+r^4/(20R^2),&r\le R,\\
2R^3/(15r),&r\ge R.
\end{cases}}\]
在$r=R$，两侧均有$\Phi=2\rho_0R^2/(15\epsilon_0)$、$E_r=2\rho_0R/(15\epsilon_0)$。图~\ref{fig:c4-v1-s}两图的球面接点均连续；没有表面电荷，法向电场没有跃变。中心有$E_r(0)=0$、$\Phi(0)=\rho_0R^2/(4\epsilon_0)$，均有限。

同一电势也满足第三讲的球对称Poisson方程：
\[\frac1{r^2}\frac{\dd}{\dd r}\left(r^2\frac{\dd\Phi}{\dd r}\right)
=-\frac{\rho(r)}{\epsilon_0}.
\]
中心正则性排除积分后的$C/r^2$电场项，球面匹配和无穷远零势再确定积分常数，所得结果与电荷积分相同。若一开始只写球外的无源方程，会漏掉全部内部源分布。

内部场强导数为
\[E_r'=\frac{\rho_0}{\epsilon_0}\left(\frac13-\frac{3r^2}{5R^2}\right).
\]
令其为零得$r^2=5R^2/9$，取区域内的正根$r=\sqrt5R/3$。此处前后$E_r'$由正变负，球外又按$r^{-2}$严格减小，故它是全空间场强最大点。将$r^2/(5R^2)=1/9$代回，
\[E_{\max}=\frac{\rho_0r}{\epsilon_0}
\left(\frac13-\frac19\right)=\frac{2\rho_0r}{9\epsilon_0},\]
所以
\[\boxed{r_{\max}=\frac{\sqrt5}{3}R,\qquad
E_{\max}=\frac{2\sqrt5}{27}\frac{\rho_0R}{\epsilon_0}}.
\]
球外场严格等于总电荷集中在球心的场，因而远势首项为$Q/(4\pi\epsilon_0r)$；在本球对称模型中，这一表达式对全部$r\ge R$都精确成立。
\begin{figure}[H]\centering
\begin{tikzpicture}[>=Stealth,line cap=round,line join=round]

\begin{scope}[x=2cm,y=7cm]
\draw[axis](0,0)--(2.65,0)node[right]{$r/R$};\draw[axis](0,-.015)--(0,.22)node[above]{$\epsilon_0E_r/(\rho_0R)$};
\draw[curve,domain=0:1,samples=81]plot(\x,{\x/3-(\x)^3/5});\draw[curve,domain=1:2.5,samples=81]plot(\x,{2/(15*(\x)^2)});
\draw[guide](1,0)--(1,.2);\node[below]at(1,0){$1$};\fill(.745356,.165635)circle(1.4pt);\node[above left]at(.745356,.165635){最大值};
\end{scope}
\begin{scope}[shift={(0,-3)},x=2cm,y=6cm]
\draw[axis](0,0)--(2.65,0)node[right]{$r/R$};\draw[axis](0,-.015)--(0,.31)node[above]{$\epsilon_0\Phi/(\rho_0R^2)$};
\draw[curve,domain=0:1,samples=81]plot(\x,{.25-(\x)^2/6+(\x)^4/20});\draw[curve,domain=1:2.5,samples=81]plot(\x,{2/(15*\x)});
\draw[guide](1,0)--(1,.27);\node[below]at(1,0){$1$};
\end{scope}

\end{tikzpicture}
\caption{同一密度决定场与势的内部形状；两图在球面$r/R=1$按无表面电荷条件匹配。}\label{fig:c4-v1-s}
\end{figure}
\end{solution}



\section{高斯积分、Gamma与Beta函数}\label{sec:c4-special}
若干连续分布与周期积分可归结为标准参数积分。变量代换先把量纲和积分范围分离，再讨论参数允许范围，不能把收敛区域外的形式等式当成积分值。

\subsection{参数微分与平方积分}
设$G_n(a)=\int_0^\infty x^n\ee^{-ax^2}\dd x$，$a>0,n>-1$。在$a$的紧正区间内，指数衰减可控制参数导数，所以
\[G_{n+2}(a)=-\frac{\dd G_n}{\dd a}.\]
换元$u=ax^2$，也可统一写成
\[\boxed{G_n(a)=\tfrac12a^{-(n+1)/2}\Gamma((n+1)/2)}.\]
参数微分把幂次提高二，不会产生奇偶序列之间的直接跳转，故需要分别知道$G_0$和$G_1$。

高斯矩首先是一个反常积分。原点附近指数趋于一，因此收敛由$x^n$决定，要求$n>-1$；无穷远处只要$a>0$，指数就压过任意固定幂次。两个端点条件各自检查后，$G_n(a)$才成为所讨论的有限函数。

对$a$求导时，被积函数变成$-x^{n+2}\ee^{-ax^2}$。若$a$在某个紧正区间内，有统一下界$a_*>0$，参数导数可由$x^{n+2}\ee^{-a_*x^2}$控制，它仍可积，因此能交换求导与积分。若同时令$a\to0^+$，这个控制会失效，不能把有限$a$下的合法求导无条件带到发散端点。

尺度变换$x=u/\sqrt a$先给出$G_n(a)=a^{-(n+1)/2}G_n(1)$。因此参数$a$的幂次由量纲缩放确定，真正需要计算的是无量纲常数$G_n(1)$。参数微分每次提高幂次二，连接偶次矩与偶次矩、奇次矩与奇次矩；$G_0,G_1$是两条递推链各自的起点。

把$G_0$平方后，原积分域是第一象限。径向高斯核只依赖$x^2+y^2$，适合换极坐标；角域却仍为$[0,\pi/2]$。几何区域决定的四分之一因子与开方时的正根选择，共同固定最后的系数。图~\ref{fig:c4-e10-q}用阴影把这一角域明确画出。

\subsection{Gamma函数与递推}
对$\Re s>0$，定义$\Gamma(s)=\int_0^\infty t^{s-1}\ee^{-t}\dd t$。在$\Gamma(s+1)$中取$u=t^s,\dd v=\ee^{-t}\dd t$作分部积分，$\dd u=st^{s-1}\dd t,v=-\ee^{-t}$，所以
\[\Gamma(s+1)=[-t^s\ee^{-t}]_0^\infty
+s\int_0^\infty t^{s-1}\ee^{-t}\dd t=s\Gamma(s).\]
$\Re s>0$保证零端边界项消失，指数保证无穷端消失。由
$\Gamma(1)=[-\ee^{-t}]_0^\infty=1$，对正整数逐次递推得$\Gamma(n)=(n-1)!$。

令$t=u^2$，则$\dd t=2u\dd u$，正半轴边界不变，故
\[\Gamma(1/2)=\int_0^\infty(u^2)^{-1/2}\ee^{-u^2}2u\dd u
=2G_0(1)=\sqrt\pi.\]
递推连接整数阶乘与半整数高斯矩，但不替代积分定义的收敛域。

对实数正参数，Gamma积分为正。半整数值可从$\Gamma(1/2)=\sqrt\pi$开始，逐次得到$\Gamma(3/2)=\sqrt\pi/2$、$\Gamma(5/2)=3\sqrt\pi/4$。这些系数常出现在球坐标径向积分中，因为平方代换把整数幂半径变成半整数幂的新变量。

\subsection{Beta函数与二维换元}
对$\Re m,\Re n>0$，定义
$B(m,n)=\int_0^1u^{m-1}(1-u)^{n-1}\dd u$。为联系Gamma函数，将乘积写成第一象限的双积分：
\[\Gamma(m)\Gamma(n)=\int_0^\infty\!\int_0^\infty
x^{m-1}y^{n-1}\ee^{-(x+y)}\dd x\dd y.\]
令$s=x+y$表示总量，$u=x/(x+y)$表示其中的比例。对$x,y>0$，恰有$s>0,0<u<1$，反变换$x=su,y=s(1-u)$一一对应，且
\[\left|\frac{\partial(x,y)}{\partial(s,u)}\right|
=\left|\begin{matrix}u&s\\1-u&-s\end{matrix}\right|=s.\]
乘积中的$x^{m-1}y^{n-1}$贡献$s^{m+n-2}u^{m-1}(1-u)^{n-1}$，再乘Jacobian的$s$，于是
\[\begin{aligned}
\Gamma(m)\Gamma(n)
&=\int_0^\infty s^{m+n-1}\ee^{-s}\dd s
\int_0^1u^{m-1}(1-u)^{n-1}\dd u\\
&=\Gamma(m+n)B(m,n).
\end{aligned}\]
故$\boxed{B(m,n)=\Gamma(m)\Gamma(n)/\Gamma(m+n)}$。参数条件保证绝对收敛，使乘积积分与换序成立。

对三角幂积分，令$u=\sin^2\theta$，则$\dd u=2\sin\theta\cos\theta\dd\theta$，$\theta=0,\pi/2$对应$u=0,1$，因此
\[\begin{aligned}
\int_0^{\pi/2}\sin^{p-1}\theta\cos^{q-1}\theta\dd\theta
&=\frac12\int_0^1u^{p/2-1}(1-u)^{q/2-1}\dd u\\
&=\frac12B(p/2,q/2),\qquad \Re p,\Re q>0.
\end{aligned}\]

$B(m,n)$在零端由$u^{m-1}$控制，在一端由$(1-u)^{n-1}$控制，故分别要求$\Re m>0,\Re n>0$。密度在端点无界仍可能可积，习题~\ref{ex:c4-h6}与~\ref{ex:c4-h16}都保留这种情况。求质心或惯量时，额外的距离幂会改变Beta的第一个参数，应根据变换后的幂次逐项辨认。

\begin{example}[高斯积分的角域与递推系数]\label{ex:c4-e10}
求$G_0(a),G_1(a),G_3(a),G_5(a)$，$a>0$。用二维积分推导$G_0$，并解释为何极角只能取$0$到$\pi/2$。 见图~\ref{fig:c4-e10-q}。
\end{example}

\begin{figure}[H]\centering
\begin{tikzpicture}[>=Stealth,line cap=round,line join=round]
\fill[gray!12](0,0)--(2.3,0)arc(0:90:2.3)--cycle;\draw[axis](0,0)--(2.7,0)node[right]{$x$};\draw[axis](0,0)--(0,2.7)node[above]{$y$};\draw[curve](2.3,0)arc(0:90:2.3);\node at(1.3,1.3){第一象限};
\end{tikzpicture}
\caption{两个半轴积分相乘，对应第一象限，不是整个平面。}\label{fig:c4-e10-q}
\end{figure}
\begin{solution}
记$G_n(a)=\int_0^\infty x^n\ee^{-ax^2}\dd x$，$a>0$。
图~\ref{fig:c4-e10-q}中，两条正半轴积分相乘对应第一象限。令$I=G_0(a)>0$，改用极坐标：
\[
 I^2=\int_0^{\pi/2}\dd\theta
 \int_0^\infty r\ee^{-ar^2}\dd r.
\]
径向积分为
\[
 \int_0^\infty r\ee^{-ar^2}\dd r
 =\left[-\frac{\ee^{-ar^2}}{2a}\right]_0^\infty
 =\frac1{2a}.
\]
故$I^2=\pi/(4a)$，取正根得
$\boxed{G_0(a)=\sqrt\pi/(2\sqrt a)}$。
若把极角误取为$0$到$2\pi$，积分区域扩大四倍，开方后的结果将多出因子二。
\begin{figure}[H]\centering
\begin{tikzpicture}[>=Stealth,line cap=round,line join=round]
\draw[axis](0,0)--(3,0)node[right]{$r$};\draw[axis](0,0)--(0,1.1)node[above]{$r\ee^{-ar^2}$};\draw[curve,domain=0:2.8]plot(\x,{\x*exp(-(\x)^2)});\node at(3.8,1.3){径向积分$1/(2a)$};\node at(3.8,.6){角域宽度$\pi/2$};\node[above]at(1.6,.5){图取$a=1$};
\end{tikzpicture}
\caption{径向与角向因子分别计算，可清楚追踪归一化系数。}\label{fig:c4-e10-s}
\end{figure}
刚才的径向积分也给出另一支种子$G_1(a)=1/(2a)$。
在任意$a>0$的小闭邻域中，指数衰减可控制求导后出现的多项式，故可在积分号下求导：
\[G_n'(a)=\int_0^\infty x^n(-x^2)\ee^{-ax^2}\dd x=-G_{n+2}(a).\]
由此
\[G_3(a)=-\frac12\frac{\dd}{\dd a}a^{-1}
=-\frac12(-a^{-2})=\boxed{\frac1{2a^2}},\]
\[G_5(a)=-\frac12\frac{\dd}{\dd a}a^{-2}
=-\frac12(-2a^{-3})=\boxed{\frac1{a^3}}.\]
偶数阶序列从$G_0$出发，奇数阶序列从图~\ref{fig:c4-e10-s}所示的径向积分$G_1$出发。半轴上的奇次幂积分非零；全实轴上对应的奇函数积分则为零，偶次幂积分是半轴的两倍。

一般地，令$u=ax^2$，则$x=\sqrt{u/a}$、$\dd x=u^{-1/2}\dd u/(2\sqrt a)$，正半轴的上下限仍为$0,\infty$。于是
\[
 G_n(a)=\frac1{2a^{(n+1)/2}}
 \int_0^\infty u^{(n-1)/2}\ee^{-u}\dd u
 =\frac{\Gamma((n+1)/2)}{2a^{(n+1)/2}}.
\]
其中$n>-1$保证原点处可积，$a>0$保证无穷远衰减。
尺度因子$a^{-(n+1)/2}$也说明$a$增大时各积分减小；整数和半整数的Gamma值分别对应上述两支递推。
\end{solution}

\section{傅里叶与拉普拉斯变换}\label{sec:c4-transforms}
频率分解适合研究各频率的线性响应，单边拉普拉斯变换则把初态直接带入方程。使用前先确定变换约定、定义域以及待求的是初值过程还是稳态频谱。

\subsection{周期信号：傅里叶级数}
周期为$T$的函数，令$\omega_0=2\pi/T$，复形式为
\[f(t)=\sum_{n\in\mathbb Z}c_n\ee^{\mathrm in\omega_0t},\quad
c_n=\frac1T\int_{t_0}^{t_0+T}f(t)\ee^{-\mathrm in\omega_0t}\dd t.
\]
归一化来自一个周期内的正交关系$T^{-1}\int\ee^{\mathrm i(n-m)\omega_0t}\dd t=\delta_{nm}$。实函数满足$c_{-n}=c_n^*$。对分段光滑函数，在连续点级数收敛到函数值，在跳跃点收敛到左右极限的平均。实三角形式采用$f=a_0/2+\sum_{n\ge1}(a_n\cos n\omega_0t+b_n\sin n\omega_0t)$，其中$c_0=a_0/2,c_n=(a_n-\mathrm ib_n)/2$。

傅里叶系数由正交投影确定。令整数$\ell=n-m\ne0$，则
\[\int_{t_0}^{t_0+T}\ee^{\ii\ell\omega_0t}\dd t
=\left[\frac{\ee^{\ii\ell\omega_0t}}{\ii\ell\omega_0}\right]_{t_0}^{t_0+T}=0,\]
因为$\ell\omega_0T=2\pi\ell$使两端指数相同；$\ell=0$时，积分为$[t]_{t_0}^{t_0+T}=T$。将有限级数乘$\ee^{-im\omega_0t}$并在一周期积分，只有$n=m$留下$Tc_m$，所以系数前必须有$1/T$；对平方可积函数可在均方极限下得到相同投影。实三角形式中，正弦余弦平方在一周期的平均为$1/2$，这固定了正整数$n$所对应的$c_n=(a_n-\ii b_n)/2$。

有限部分和是光滑周期函数，不能在整个区间上一致逼近跳跃函数；在满足相应分段光滑条件时，它可在连续点逐点收敛，在跳跃点收敛到左右平均。均方收敛考察的是误差平方的积分，允许误差集中在越来越窄的区域。因此方波的局部过冲与Parseval关系中的均方误差趋零可以同时存在。

若目标是解线性周期驱动问题，可把每个谐波分别代入频率响应，再按系数叠加；线性保证不同频率不在方程中互相相乘。非线性模型则可能产生新频率，不能直接沿用同样的逐谐波独立求解。傅里叶展开的用途取决于随后所用动力学的结构。

\subsection{非周期连续信号：傅里叶变换}
本讲约定
\[F(\omega)=\int_{-\infty}^\infty f(t)\ee^{-\mathrm i\omega t}\dd t,\quad
f(t)=\frac1{2\pi}\int_{-\infty}^\infty F(\omega)\ee^{\mathrm i\omega t}\dd\omega.
\]
时间平移$f(t-t_0)$对应频域乘$\ee^{-\mathrm i\omega t_0}$；在边界衰减充分时，时间微分对应乘$\mathrm i\omega$。系数$2\pi$与角频率约定必须配套，不能把以普通频率为变量的公式直接混入。

连续傅里叶变换的频率变量可取角频率$\omega$或普通频率$\nu$，二者满足$\omega=2\pi\nu$。本讲使用角频率，逆式前的$1/(2\pi)$与$\dd\omega$配对；换成$\ee^{-2\pi i\nu t}$约定时，逆式的微元也变成$\dd\nu$。归一化差别来自变量替换，而非两种物理频谱。

普通绝对收敛积分定义下，若$f\in L^1$，变换$F$存在且连续；要直接使用绝对收敛的逆积分，可再要求$F\in L^1$。更一般的平方可积函数在均方意义下也能定义变换，跳跃或不绝对可积情形则需说明所用的极限意义。本讲的矩形脉冲与高斯函数具有明确的积分计算，可在各自条件下讨论。

平移性质可从$t=u+t_0$直接推出，微分及全轴边界保持不变，指数分出$\ee^{-\ii\omega t_0}$，即得前面的相位因子。时间伸缩则令$t=au$，$a>0$时$\dd t=a\dd u$，所以
\[\int_{-\infty}^{\infty}f(t/a)\ee^{-\ii\omega t}\dd t
=a\int_{-\infty}^{\infty}f(u)\ee^{-\ii(a\omega)u}\dd u=aF(a\omega).\]
时间微分的性质来自分部积分：
\[\int_{-\infty}^{\infty}f'(t)\ee^{-\ii\omega t}\dd t
=[f(t)\ee^{-\ii\omega t}]_{-\infty}^{\infty}
+\ii\omega F(\omega)=\ii\omega F(\omega),\]
其中端点项消失需有相应衰减，积分也需存在。图~\ref{fig:c4-h8-q}与~\ref{fig:c4-h8-s}分别展示脉冲和频谱：脉宽减小会使谱宽增大，同时改变零频面积；两个横轴应分别使用时间和角频率的单位。

\subsection{单边拉普拉斯变换与初值}
对$t\ge0$的分段连续函数，定义$F(s)=\mathcal L[f]=\int_0^\infty f(t)\ee^{-st}\dd t$。若$f$的绝对值不快于$\ee^{\sigma_0t}$增长，取$\Re s>\sigma_0$可保证无穷端收敛。在函数及所需导数具备相应正则性和增长条件时，分部积分得到
\[\mathcal L[f']=sF-f(0^+),\quad
\mathcal L[f'']=s^2F-sf(0^+)-f'(0^+).
\]
因此微分方程化为代数方程时，初值以明确的附加项出现。常用对为$1\leftrightarrow1/s$、$\ee^{at}\leftrightarrow1/(s-a)$、$\sin\omega t\leftrightarrow\omega/(s^2+\omega^2)$、$\cos\omega t\leftrightarrow s/(s^2+\omega^2)$。卷积$(f*g)(t)=\int_0^tf(\tau)g(t-\tau)\dd\tau$对应乘积$F(s)G(s)$，这与上一讲因果响应的叠加积分相接。

\Needspace{8\baselineskip}
单边Laplace变换把初值纳入代数式，其导数性质来自带边界项的分部积分：
\[\int_0^\infty f'(t)\ee^{-st}\dd t
=[f(t)\ee^{-st}]_0^\infty+sF(s)=sF(s)-f(0^+).\]
无穷端项消失需要$s$位于相应收敛半平面；零端项则正是初值。再对$f'$应用同一公式，
\[\mathcal L[f'']=s\mathcal L[f']-f'(0^+)
=s(sF-f(0^+))-f'(0^+),\]
即得前面的二阶导数式，初位移和初速度分别留下不同的$s$幂次项。

常用变换对可由指数积分建立。在$\Re(s-a)>0$时，
\[\mathcal L[\ee^{at}]=\int_0^\infty\ee^{-(s-a)t}\dd t
=\left[-\frac{\ee^{-(s-a)t}}{s-a}\right]_0^\infty=\frac1{s-a}.\]
取$a=0$得到常数的变换；取$a=\pm\ii\omega$再用Euler公式，
\[\mathcal L[\sin\omega t]=\frac1{2\ii}
\left(\frac1{s-\ii\omega}-\frac1{s+\ii\omega}\right)
=\frac\omega{s^2+\omega^2},\]
\[\mathcal L[\cos\omega t]=\frac12
\left(\frac1{s-\ii\omega}+\frac1{s+\ii\omega}\right)
=\frac s{s^2+\omega^2}.\]

斜坡输入的变换也由分部积分求出。在$\Re s>0$时，
\[
 \mathcal L[t]=\left[-\frac{t\ee^{-st}}s\right]_0^\infty
 +\frac1s\int_0^\infty\ee^{-st}\dd t
 =\frac1{s^2}.
\]
再令$\tau>0$，延迟后的因果函数为$H(t-\tau)f(t-\tau)$。取$u=t-\tau$，积分下限由$\tau$移到零，得到
\[
\begin{aligned}
 \mathcal L[H(t-\tau)f(t-\tau)]
 &=\int_\tau^\infty f(t-\tau)\ee^{-st}\dd t\\
 &=\ee^{-s\tau}\int_0^\infty f(u)\ee^{-su}\dd u
 =\ee^{-s\tau}F(s).
\end{aligned}
\]
阶跃因子使延迟前的函数为零；仅把$t$改成$t-\tau$而不截去此前区间，不能得到这一因果时移关系。

这里的$0^+$约定适用于先给定正侧初态、再求$t>0$的普通方程。若$t=0$另有冲量，应先由跨零积分确定初态跳变，再将得到的正侧初值代入；若直接把冲量放入分布变换，则必须采用与负侧初态一致的另一套边界约定，不能把两种处理重复相加。

变换的代数表达式还附有收敛域。$1/(s-a)$来自$\ee^{at}$时，原积分要求$\Re s>\Re a$，不能只保留一个分式而忘记其因果性与增长条件。解方程得到$X(s)$后，部分分式、已知变换对与卷积可用于反演；每一项的初值与长时行为还应回到原方程解释。

\begin{example}[方波级数与有限谐波近似]\label{ex:c4-e11}
周期$2\pi$的方波在$0<t<\pi$取$1$，$-\pi<t<0$取$-1$。求实傅里叶级数，说明跳跃点的取值与有限项近似的局限。 见图~\ref{fig:c4-e11-q}。
\end{example}

\begin{figure}[H]\centering
\begin{tikzpicture}[>=Stealth,line cap=round,line join=round]
\draw[axis](-3.65,0)--(3.65,0)node[right]{$t$};\draw[axis](0,-1.5)--(0,1.6)node[above]{$f$};
\draw[curve](-3.14159,-1)--(0,-1) (0,1)--(3.14159,1);
\foreach\x/\y in{-3.14159/-1,0/-1,0/1,3.14159/1}{\draw[fill=white](\x,\y)circle(1.6pt);}
\node[below]at(-3.14159,0){$-\pi$};\node[below]at(3.14159,0){$\pi$};
\node[left]at(0,1){$1$};\node[right]at(0,-1){$-1$};
\end{tikzpicture}
\caption{一个周期内的方波，端点用空心点表示区间端值未指定；跳跃处的竖直连接段不属于函数图像。}\label{fig:c4-e11-q}
\end{figure}
\begin{solution}
方波为奇函数，直流项与余弦系数均为零。图~\ref{fig:c4-e11-q}还满足半周期反号，偶数次正弦谐波也随之消失。下面分别用这两种对称简化系数积分。


由奇性，$a_0=a_n=0$。正弦系数为
\[b_n=\frac2\pi\int_0^\pi\sin nt\dd t
=\frac2\pi\left[-\frac{\cos nt}{n}\right]_0^\pi
=\frac{2[1-(-1)^n]}{\pi n}.
\]
于是
\[\boxed{f(t)=\frac4\pi\sum_{j=0}^\infty\frac{\sin[(2j+1)t]}{2j+1}}\]
在各连续点成立，在$t=k\pi$跳跃处收敛到零。有限项是连续函数，不能在跳跃点附近一致逼近不连续方波。图~\ref{fig:c4-e11-s}保留前三个非零谐波，跳跃附近出现过冲；增加项数使过渡区收窄，但不能消除极限中的相对过冲高度。

若把跳跃点单独定义为$1$，积分系数也不改变，因为有限个点不影响积分，级数在该点仍取左右极限平均。傅里叶展开恢复的是几乎处处的函数信息，不记录任意指定的孤立点值。


图~\ref{fig:c4-e11-s}所用的前三项为
\[S_5(t)=\frac4\pi\left(\sin t+\frac{\sin3t}{3}+\frac{\sin5t}{5}\right),\]
它保持奇对称性，且在$t=0,\pi$处均取零，与左右极限平均值一致。

\begin{figure}[H]\centering
\begin{tikzpicture}[>=Stealth,line cap=round,line join=round]
\begin{scope}[x=0.80000cm,y=0.80000cm]\draw[axis](-3.14159,0)--(3.39159,0)node[right]{$t$};\draw[axis](0,-1.6)--(0,1.6)node[above]{$S_5$};\draw[curve,domain=-3.14159:3.14159,samples=101]plot(\x,{4/pi*(sin(\x r)+sin(3*\x r)/3+sin(5*\x r)/5)});\draw[dashed](-3.14159,-1)--(0,-1) (0,1)--(3.14159,1);\end{scope}
\end{tikzpicture}
\caption{前三个奇次谐波的部分和；跳跃处趋向左右平均值。}\label{fig:c4-e11-s}
\end{figure}
\end{solution}

\Needspace{8\baselineskip}
\begin{example}[用拉普拉斯变换处理振子初值]\label{ex:c4-e12}
求$\ddot x+2\gamma\dot x+\omega_0^2x=f_0$在$x(0)=0,\dot x(0)=v_0$下的解，取$0<\gamma<\omega_0$。再说明任意时变右端$f(t)$如何进入答案。 见图~\ref{fig:c4-e12-q}。
\end{example}

\begin{figure}[H]\centering
\begin{tikzpicture}[>=Stealth,line cap=round,line join=round]
\node[draw,align=center,inner sep=7pt](A)at(0,1.7){$\ddot x+2\gamma\dot x+\omega_0^2x=f_0$\\$x(0)=0,\quad\dot x(0)=v_0$};\node[draw,inner sep=7pt](B)at(0,-.4){$(s^2+2\gamma s+\omega_0^2)X=v_0+f_0/s$};\draw[axis](A)--(B)node[midway,right]{$\mathcal L$};
\end{tikzpicture}
\caption{初速度与恒外力分别产生代数式右端的不同项。}\label{fig:c4-e12-q}
\end{figure}
\begin{solution}
图~\ref{fig:c4-e12-q}中的恒驱动从$t=0$开启。作单边拉普拉斯变换后，驱动产生$f_0/s$，初速度通过二阶导数的零端边界项进入代数式。以下分别反演这两部分。


取单边变换，二阶导给出$s^2X-v_0$，一阶导因$x(0)=0$成为$sX$。所以
\[(s^2X-v_0)+2\gamma sX+\omega_0^2X=\frac{f_0}{s}.\]
将二次项配方为$(s+\gamma)^2+\omega_0^2-\gamma^2$，得
\[X(s)=\frac{v_0}{(s+\gamma)^2+\omega_d^2}
+\frac{f_0}{s[(s+\gamma)^2+\omega_d^2]},\quad\omega_d=\sqrt{\omega_0^2-\gamma^2}.
\]
由正弦的变换对及指数平移性质，定义
\[G(t)=\frac{\ee^{-\gamma t}\sin\omega_dt}{\omega_d},\qquad
\mathcal L[G]=\frac1{(s+\gamma)^2+\omega_d^2}.\]
第一项反演为$v_0G(t)$；第二项为$f_0\int_0^tG(u)\dd u$。为算此积分，记
\begin{align*}
J_s&=\int\ee^{-\gamma u}\sin\omega_du\dd u,\\
J_c&=\int\ee^{-\gamma u}\cos\omega_du\dd u.
\end{align*}
分别对三角因子积分，分部积分给出
\[J_s=-\frac{\ee^{-\gamma u}\cos\omega_du}{\omega_d}
-\frac\gamma{\omega_d}J_c,\qquad
J_c=\frac{\ee^{-\gamma u}\sin\omega_du}{\omega_d}
+\frac\gamma{\omega_d}J_s.\]
代入并移项，利用$\omega_d^2+\gamma^2=\omega_0^2$，有
\[\int G(u)\dd u
=-\frac{\ee^{-\gamma u}}{\omega_0^2}
\left(\cos\omega_du+\frac\gamma{\omega_d}\sin\omega_du\right)+C.\]
取$0,t$两端，得到
\[\int_0^tG(u)\dd u=\frac{1-\ee^{-\gamma t}[\cos\omega_dt+(\gamma/\omega_d)\sin\omega_dt]}{\omega_0^2}.\]
因此
\[\boxed{x(t)=\frac{v_0}{\omega_d}\ee^{-\gamma t}\sin\omega_dt
+\frac{f_0}{\omega_0^2}\left[1-\ee^{-\gamma t}\left(\cos\omega_dt+\frac\gamma{\omega_d}\sin\omega_dt\right)\right]}.
\]
图~\ref{fig:c4-e12-s}中解由初态瞬变与外源响应相加，长时趋于静平衡$f_0/\omega_0^2$。初态可从响应核的导数核对：
\[G'=\ee^{-\gamma t}\left(\cos\omega_dt-\frac\gamma{\omega_d}\sin\omega_dt\right),\]
\[G''=\ee^{-\gamma t}\left[-2\gamma\cos\omega_dt
+\frac{\gamma^2-\omega_d^2}{\omega_d}\sin\omega_dt\right].\]
两次乘积求导分别产生指数因子和三角因子的导数，故$G(0)=0$、$G'(0)=1$、$G''(0)=-2\gamma$。由
$x'=v_0G'+f_0G$、$x''=v_0G''+f_0G'$，得$x(0)=0$、$x'(0)=v_0$、$x''(0)=f_0-2\gamma v_0$，与原方程一致。

任意右端$f(t)$只需把$f_0/s$换成$F(s)$，于是
\[x(t)=v_0G(t)+\int_0^tG(t-\tau)f(\tau)\dd\tau.
\]
若原方程右侧是物理力$F_{\rm ext}(t)$，这里$f=F_{\rm ext}/m$，不可遗漏质量。拉普拉斯变换给出的因果卷积与直接的响应函数法相同。



\begin{figure}[H]\centering
\begin{tikzpicture}[>=Stealth,line cap=round,line join=round]
\begin{scope}[x=1.00000cm,y=1.05882cm]\draw[axis](0,0)--(6.25,0)node[right]{$\omega_0t$};\draw[axis](0,0)--(0,1.7)node[above]{$\omega_0^2x/f_0$};\draw[curve,domain=0:6,samples=101]plot(\x,{1-exp(-.3*\x)*(cos(.953939*\x r)+.314485*sin(.953939*\x r))+.4*exp(-.3*\x)*sin(.953939*\x r)/.953939});\draw[guide](0,1)--(6,1);\end{scope}
\end{tikzpicture}
\caption{取$\gamma/\omega_0=.3,\ v_0\omega_0/f_0=.4$；长时趋于新的静平衡。}\label{fig:c4-e12-s}
\end{figure}
\end{solution}

\subsection{同一有限输入的分段、卷积与变换}
对例~\ref{ex:c2-v1}的有限斜坡，分段初值解与卷积解已得到相同结果。沿用该题的零初态，令$H(t)$为单位阶跃，则在$t=\tau$撤去的斜坡为
\[F(t)=\frac{F_0}{\tau}\{tH(t)-(t-\tau)H(t-\tau)-\tau H(t-\tau)\}.
\]
最后的延迟常值项把$\tau$之后的$F_0$也减去；若只作前两个斜坡之差，输入会停在$F_0$。图~\ref{fig:c4-ramp-decomposition}把开启斜坡与两个延迟修正画在同一时间轴上。

\begin{figure}[H]\centering
\begin{tikzpicture}[>=Stealth,line cap=round,line join=round]

\draw[axis](-.2,0)--(5.7,0)node[right]{$t/\tau$};\draw[axis](0,-.25)--(0,2.9)node[above]{$F/F_0$};
\draw[curve](0,0)--(1.5,1.3)--(1.5,0)--(4.5,0);
\draw[dashed](0,0)--(1.5,1.3)--(3,2.6);
\node[below]at(1.5,0){$1$};\node[left]at(0,1.3){$1$};\draw[guide](0,1.3)--(1.5,1.3);\node[right,align=left]at(3.3,1.1){虚线：未截断斜坡\\[5pt]实线：有限斜坡};

\end{tikzpicture}
\caption{有限斜坡须同时减去延迟斜坡与延迟常值，才会在撤力后归零。}\label{fig:c4-ramp-decomposition}
\end{figure}

使用$\mathcal L[t]=1/s^2$以及延迟性质，有
\[\widehat F(s)=\frac{F_0}{\tau}
\left(\frac{1-\ee^{-s\tau}}{s^2}-\frac{\tau\ee^{-s\tau}}s\right),
\qquad X(s)=\frac{\widehat F(s)}{m(s^2+\omega^2)}.\]
对无阻尼核$G_{\rm free}(t)=\sin\omega t/\omega$，先后累积一次和两次，得到
\[S(t)=\int_0^tG_{\rm free}(u)\dd u
=\left[-\frac{\cos\omega u}{\omega^2}\right]_0^t
=\frac{1-\cos\omega t}{\omega^2},\]
\[R(t)=\int_0^tS(u)\dd u
=\frac1{\omega^2}\left[u-\frac{\sin\omega u}{\omega}\right]_0^t
=\frac{t-\sin\omega t/\omega}{\omega^2}.\]
把这两个响应在负时间定义为零。它们分别是$1/[s(s^2+\omega^2)]$与$1/[s^2(s^2+\omega^2)]$的反变换，且$R'=S$。故
\[\boxed{x(t)=\frac{F_0}{m\tau}
[R(t)-R(t-\tau)-\tau S(t-\tau)]}.\]
$0<t<\tau$时只剩第一项，即第二讲的受力段；$t>\tau$时展开三角函数，恢复以$x_\tau,v_\tau$为初值的自由运动。分段法直接使用连接状态，卷积法把输入限制在$0<s<\min(t,\tau)$，变换法把同一边界写成$\ee^{-s\tau}$；三者处理的是同一个因果过程。

由$R(0)=R'(0)=S(0)=S'(0)=0$，统一表达式在$\tau$处的位置与速度连续。$S''(0^+)=1$，因而加速度跳变为$-F_0/m$，与原方程一致。这里$S$始终表示阶跃响应，$H$只表示阶跃输入函数，二者的物理量纲和用途不同。



\section{变换中的守恒关系与收敛条件}
变换后的代数表达还须对应一个实际存在的时域或空间域函数。以下用能量关系、收敛区域和极点位置解释变换结果，避免只完成形式运算。

\subsection{正交投影与Parseval关系}
在周期$T$上，复指数是正交基。若先考虑有限和$f_N=\sum_{|n|\le N}c_n\ee^{\mathrm in\omega_0t}$，把$|f_N|^2$展开，交叉项的周期积分全部消失，因此
\[\begin{aligned}
\frac1T\int_{t_0}^{t_0+T}|f_N|^2\dd t
&=\sum_{|n|,|m|\le N}c_nc_m^*
\frac1T\int_{t_0}^{t_0+T}\ee^{\ii(n-m)\omega_0t}\dd t\\
&=\sum_{|n|,|m|\le N}c_nc_m^*\delta_{nm}
=\sum_{|n|\le N}|c_n|^2.
\end{aligned}\]
对平方可积函数，以均方极限过渡得到Parseval关系。实三角系数的形式为
\[\frac1T\int|f|^2\dd t=\frac{a_0^2}{4}+\frac12\sum_{n\ge1}(a_n^2+b_n^2).
\]
它把时域均方量与频域各谐波贡献相连。部分和未包含的平方系数之和，就是截断后的均方误差；这种误差描述与跳跃点附近的最大点态误差不同。

连续傅里叶变换中的归一化可从逆式推得。先取函数及所需导数均快速衰减的情形，使下面的积分可绝对换序。将$f$的逆变换代入$f\overline f$，有
\[
\begin{aligned}
 \int_{-\infty}^{\infty}|f(t)|^2\dd t
 &=\frac1{2\pi}\int_{-\infty}^{\infty}F(\omega)
 \left[\int_{-\infty}^{\infty}\overline{f(t)}\ee^{\ii\omega t}\dd t\right]\dd\omega\\
 &=\frac1{2\pi}\int_{-\infty}^{\infty}|F(\omega)|^2\dd\omega.
\end{aligned}
\]
方括号是$\overline{F(\omega)}$。对一般平方可积函数，可由快速衰减函数的均方逼近延拓这个等式。

有限DFT则利用习题~\ref{ex:c4-h10}中由几何级数证明的正交和
\[
 \sum_{k=0}^{N-1}\ee^{2\pi\ii k(n-m)/N}=N\delta_{nm},
 \qquad 0\le n,m<N.
\]
将$X_k$与其共轭的有限和相乘，再对$k$求和：
\[
\begin{aligned}
 \sum_{k=0}^{N-1}|X_k|^2
 &=\sum_{n,m=0}^{N-1}x_n\overline{x_m}
 \sum_{k=0}^{N-1}\ee^{-2\pi\ii k(n-m)/N}\\
 &=N\sum_{n=0}^{N-1}|x_n|^2.
\end{aligned}
\]
因此时域平方和等于频域平方和的$1/N$。连续式的$1/(2\pi)$与离散式的$1/N$各由所用变换的正交关系确定。

\subsection{卷积定理的积分域}
因果卷积$h(t)=\int_0^tf(\tau)g(t-\tau)\dd\tau$的拉普拉斯变换，对区域$0\le\tau\le t<\infty$换元$u=t-\tau$，成为第一象限$\tau,u\ge0$上的积分：
\[H(s)=\int_0^\infty\!\int_0^\infty f(\tau)g(u)\ee^{-s(\tau+u)}\dd u\dd\tau=F(s)G(s).
\]
交换积分的步骤需要绝对收敛或其他可适用的定理条件。若两个变换的收敛域分别为$\Re s>\sigma_f,\Re s>\sigma_g$，通常至少在其公共右半平面里推导成立。代数式的解析延拓可能更大，但延拓后的每一点未必仍由原积分收敛定义。

\subsection{稳定性与变换域的极点}
线性系统若响应中含$\ee^{pt}$，则极点$\Re p<0$对应衰减，$\Re p>0$对应增长；虚轴上的非零单极点对应不衰减振荡，原点单极点对应常量。重极点可能带时间多项式并导致增长。对实系数系统，复极点成共轭对，保证实响应。

由拉普拉斯代数式反演时，要保留初值项和因果性。终值关系
\[\lim_{t\to\infty}f(t)=\lim_{s\to0}sF(s)\]
需要相应的收敛条件。若响应有持续振荡、增长模态或不收敛项，右侧代数极限可能存在而时域终值不存在。应先由极点或直接时间解确认收敛，再用其求长时值。

对普通因果函数的真有理Laplace变换$F(s)$，一个常用的充分条件是：约去可约因子后，$sF(s)$的全部极点严格位于左半平面。此时$F(s)$在原点至多有一个单极点。记$L=\lim_{s\to0}sF(s)$，作部分分式分解可写成
\[
 F(s)=\frac Ls+F_d(s),
\]
其中$F_d$的极点均在严格左半平面。第一项反演为常量$L$，其余项只产生衰减指数乘有限次时间多项式，所以$f(t)\to L$。若有非零虚轴极点或右半平面极点，这个论证便不能使用；习题~\ref{ex:c4-h15}给出代数极限存在而时间终值不存在的例子。应先核实极点条件，再取终值极限。

习题~\ref{ex:c4-h13}用方波的单位均方值确定系数平方和，再由截断后遗漏的系数计算均方误差。


\section{延伸讨论：离散采样与频谱}
连续模型变成离散数据后，不同频率可能产生相同样本。先规定采样间隔和观测长度，再区分混叠、频率分辨率与数值误差。

对绝对可和序列$x_n$，定义$X(\omega)=\sum_{n\in\mathbb Z}x_n\ee^{-\mathrm in\omega}$。它是$2\pi$周期函数，逆变换只能在任意一个完整周期上积分：
\[\boxed{x_n=\frac1{2\pi}\int_{-\pi}^\pi X(\omega)\ee^{\mathrm in\omega}\dd\omega}.
\]
将周期频谱在整个实轴积分，通常会重复无穷次而不收敛。代入级数并交换求和积分后，利用$\int_{-\pi}^\pi\ee^{\mathrm i(n-m)\omega}\dd\omega=2\pi\delta_{nm}$即可证明逆式。

长度$N$的有限序列DFT则是$X_k=\sum_{n=0}^{N-1}x_n\ee^{-2\pi\mathrm ikn/N}$，逆式$x_n=N^{-1}\sum_{k=0}^{N-1}X_k\ee^{2\pi\mathrm ikn/N}$。它把有限数据视为周期延拓，频率只取$N$个离散点。DTFT有连续频率变量，DFT有有限频率样点；采样与截断还会带来混叠与谱泄漏，不能把两者统称成同一个积分变换。

DTFT的周期性来自整数样点：$\ee^{-in(\omega+2\pi)}=\ee^{-in\omega}$。这里$\omega$是每个样点的无量纲角频率；若采样间隔为$\Delta t$，连续物理角频率为$\omega/\Delta t$。周期频谱在全实轴重复，逆变换选择一段长度$2\pi$的区间恰好使整数指数正交。

DFT则先截取$N$个样点，再只在$\omega_k=2\pi k/N$上求频谱值。它隐含对这$N$个样点作周期延拓，因此频率格点间隔对应物理值$2\pi/(N\Delta t)$。增加真实观测时长能细化频率分辨；仅在末尾补零虽然增加绘图样点，却不增加原记录的观测信息。

混叠与泄漏的原因不同。相差$2\pi/\Delta t$的连续复指数在所有采样点上完全相同，属于采样造成的混叠；不在DFT格点上的频率被有限窗口截断后，会分布到多个格点，属于泄漏。习题~\ref{ex:c4-h10}选取整格点余弦，先展示两个精确谱点，再说明采样频率的等价类。




\section{习题与解答}

\Needspace{8\baselineskip}
\begin{exercise}[换序后的抛物线区域]\label{ex:c4-h1}
计算$I=\int_0^1\dd x\int_{x^2}^{x}y\dd y$，并交换积分次序复算。 见图~\ref{fig:c4-h1-q}。
\end{exercise}

\begin{figure}[H]\centering
\begin{tikzpicture}[>=Stealth,line cap=round,line join=round]
\draw[axis](0,0)--(2.4,0)node[right]{$x$};\draw[axis](0,0)--(0,2.4)node[above]{$y$};\draw[curve](0,0)--(2,2);\draw[curve,domain=0:2]plot(\x,{(\x)^2/2});\draw[thick](1,.5)--(1,1);
\end{tikzpicture}
\caption{竖直扫描在每个$x$上使用$x^2\le y\le x$。}\label{fig:c4-h1-q}
\end{figure}
\begin{solution}
图~\ref{fig:c4-h1-q}中的竖直细条满足
$0\le x\le1$、$x^2\le y\le x$。按原次序，
\[
 I=\int_0^1\left[\frac{y^2}{2}\right]_{x^2}^x\dd x
 =\frac12\left[\frac{x^3}{3}-\frac{x^5}{5}\right]_0^1
 =\frac1{15}.
\]
\begin{figure}[H]\centering
\begin{tikzpicture}[>=Stealth,line cap=round,line join=round]
\draw[axis](0,0)--(2.4,0)node[right]{$x$};\draw[axis](0,0)--(0,2.4)node[above]{$y$};\draw[curve](0,0)--(2,2);\draw[curve,domain=0:2]plot(\x,{(\x)^2/2});\draw[thick](.8,.8)--(1.264911,.8);\node[below]at(.8,.8){$y$};\node[right]at(1.3,.8){$\sqrt y$};
\end{tikzpicture}
\caption{水平扫描由直线进入、从抛物线离开。}\label{fig:c4-h1-s}
\end{figure}
改用图~\ref{fig:c4-h1-s}中的水平细条。由$x^2\le y$及$x\ge0$得$x\le\sqrt y$，由$y\le x$得$x\ge y$。
整个区域的$y$范围为$[0,1]$，且该区间上$y\le\sqrt y$，所以左端在直线$x=y$，右端在抛物线$x=\sqrt y$。换序后
\[
 \begin{aligned}
 I&=\int_0^1\dd y\int_y^{\sqrt y}y\dd x
 =\int_0^1(y^{3/2}-y^2)\dd y\\
 &=\left[\frac25y^{5/2}-\frac13y^3\right]_0^1
 =\boxed{\frac1{15}}.
 \end{aligned}
\]
两种细条各覆盖同一区域一次。交换积分次序须重新确定内外边界，不能只交换积分符号而沿用旧积分限。
\end{solution}

\Needspace{8\baselineskip}
\begin{exercise}[圆环薄片的转动惯量]\label{ex:c4-h2}
半径$R$的薄圆盘面密度$\sigma=\sigma_0(1+\alpha r^2/R^2)$，$\alpha>-1$。求质量、绕中心法线的转动惯量与回转半径。 见图~\ref{fig:c4-h2-q}。
\end{exercise}

\begin{figure}[H]\centering
\begin{tikzpicture}[>=Stealth,line cap=round,line join=round]
\draw[curve](0,0)circle(1.5);\draw[curve](0,0)circle(.9);\draw[axis](0,0)--(.9,0)node[midway,above]{$r$};\node[right]at(1.6,.8){$\sigma(r)$};
\end{tikzpicture}
\caption{面密度随半径变化，圆环内仍保持转动对称。}\label{fig:c4-h2-q}
\end{figure}
\begin{solution}
图~\ref{fig:c4-h2-q}中的同心圆环上密度相同，可先合并周向积分。取图~\ref{fig:c4-h2-s}的窄环，质量微元为
\[\dd m=2\pi\sigma_0(1+\alpha r^2/R^2)r\dd r.\]
质量与轴向惯量分别累计$\dd m$和$r^2\dd m$，所以
\[\begin{aligned}
M&=2\pi\sigma_0\int_0^R\left(r+\frac{\alpha r^3}{R^2}\right)\dd r\\
&=2\pi\sigma_0\left[\frac{r^2}{2}+\frac{\alpha r^4}{4R^2}\right]_0^R
=\boxed{\pi\sigma_0R^2(1+\alpha/2)},\\
I_z&=2\pi\sigma_0\int_0^R\left(r^3+\frac{\alpha r^5}{R^2}\right)\dd r\\
&=2\pi\sigma_0\left[\frac{r^4}{4}+\frac{\alpha r^6}{6R^2}\right]_0^R
=\boxed{\pi\sigma_0R^4(1/2+\alpha/3)}.
\end{aligned}\]
回转半径定义为$k_g=\sqrt{I_z/M}$，故
\[\boxed{k_g=R\sqrt{\frac{1/2+\alpha/3}{1+\alpha/2}}}.\]
$\alpha=0$恢复均匀圆盘的$R/\sqrt2$。为确定密度变化的影响，将比值化为
\[\frac{k_g^2}{R^2}=\frac{3+2\alpha}{3(2+\alpha)},\qquad
\frac{\dd}{\dd\alpha}\frac{k_g^2}{R^2}
=\frac{2(2+\alpha)-(3+2\alpha)}{3(2+\alpha)^2}
=\frac1{3(2+\alpha)^2}>0.\]
外缘密度增强时，远离转轴的质量权重增加，回转半径随之增大；$\alpha\to\infty$时趋于$R\sqrt{2/3}$。非均匀圆盘的惯量必须按实际分布积分，不能只在均匀圆盘公式中替换总质量。

\begin{figure}[H]\centering
\begin{tikzpicture}[>=Stealth,line cap=round,line join=round]
\draw[axis](0,-1.8)--(0,1.8)node[above]{$z$};\draw[curve](0,0)ellipse(1.6 and .6);\draw[curve](0,0)ellipse(1 and .375);\draw[axis](0,0)--(1,0)node[midway,above]{$r$};\node at(3,.7){$\dd I_z=r^2\dd m$};
\end{tikzpicture}
\caption{每一圆环绕轴距离相同，先形成质量元，再乘$r^2$。}\label{fig:c4-h2-s}
\end{figure}
\end{solution}

\Needspace{8\baselineskip}
\begin{exercise}[椭球体积与惯量]\label{ex:c4-h3}
均匀实心椭球$x^2/a^2+y^2/b^2+z^2/c^2\le1$，密度$\rho$。求体积和绕$z$轴的转动惯量。 见图~\ref{fig:c4-h3-q}。
\end{exercise}

\begin{figure}[H]\centering
\begin{tikzpicture}[>=Stealth,line cap=round,line join=round]
\draw[curve](0,0)ellipse(2 and 1.2);\draw[dashed](0,0)ellipse(2 and .4);\draw[axis](0,-1.5)--(0,1.7)node[above]{$z$};\draw[axis](0,0)--(2,0)node[midway,below]{$a$};
\end{tikzpicture}
\caption{把椭球伸缩成单位球时，体积元也发生伸缩。}\label{fig:c4-h3-q}
\end{figure}
\begin{solution}
椭球通过坐标伸缩映为单位球，但体积元必须一同伸缩。图~\ref{fig:c4-h3-q}中三个半轴分别给三个独立比例，合成Jacobian$abc$。惯量中的$x^2+y^2$也要分别带上$a^2,b^2$，不能仅在球体公式中替换一个“平均半径”。


取$x=au,y=bv,z=cw$，参数域为单位球$B$。Jacobian为
\[\left|\frac{\partial(x,y,z)}{\partial(u,v,w)}\right|
=\left|\det\begin{pmatrix}a&0&0\\0&b&0\\0&0&c\end{pmatrix}\right|=abc.\]
由单位球径向积分，
\[V=abc\,4\pi\int_0^1r^2\dd r
=4\pi abc\left[\frac{r^3}{3}\right]_0^1=\frac{4\pi abc}{3},\qquad M=\rho V.\]
图~\ref{fig:c4-h3-s}中伸缩后的惯量为
\[I_z=\rho abc\iiint_{u^2+v^2+w^2\le1}(a^2u^2+b^2v^2)\dd u\dd v\dd w.
\]
由球对称性，三个平方积分相同，其和为
\[\iiint_B(u^2+v^2+w^2)\dd V
=4\pi\int_0^1r^4\dd r
=4\pi\left[\frac{r^5}{5}\right]_0^1=\frac{4\pi}5.\]
旋转对称性使三个坐标平方各占三分之一，所以每一项为$4\pi/15$。代回图中的伸缩变换，
\[I_z=\rho abc(a^2+b^2)\frac{4\pi}{15},\quad
M=\rho abc\frac{4\pi}3.\]
两式相除，$I_z/M=(a^2+b^2)(4\pi/15)/(4\pi/3)$，得到
\[\boxed{I_z=\frac M5(a^2+b^2)}.\]
变换后的域虽然是球，目标函数仍分别包含$a^2,b^2$，三个方向的伸缩不能合并成一个长度比例。绕$x,y$轴的结果可循环交换半轴取得。

图~\ref{fig:c4-h3-s}的伸缩保留了关于各坐标面的反射对称，故$\int xy\dd m$等惯量积均为零，坐标轴就是主轴。令$a=b=c=R$，式子恢复实心球的$2MR^2/5$。

\begin{figure}[H]\centering
\begin{tikzpicture}[>=Stealth,line cap=round,line join=round]
\draw[curve](0,0)ellipse(1.6 and .9);\draw[axis](2,0)--(3.4,0)node[midway,above]{伸缩};\draw[curve](5,0)circle(1);\node at(5,-1.4){$u^2+v^2+w^2\le1$};\node at(0,-1.4){$\dd V=abc\dd u\dd v\dd w$};
\end{tikzpicture}
\caption{区域被化为单位球，Jacobian和距离平方都须保留。}\label{fig:c4-h3-s}
\end{figure}
\end{solution}

\Needspace{8\baselineskip}
\begin{exercise}[曲面通量的定向]\label{ex:c4-h4}
求$\vct A=(0,0,z)$通过曲面$z=k(x^2+y^2),r\le R$的向上通量，取$k,R>0$，再用封闭曲面验证。 见图~\ref{fig:c4-h4-q}。
\end{exercise}

\begin{figure}[H]\centering
\begin{tikzpicture}[>=Stealth,line cap=round,line join=round]
\draw[curve,domain=-1.5:1.5]plot(\x,{.6*(\x)^2});\draw[dashed](-1.5,1.35)--(1.5,1.35);\draw[axis](0,-.2)--(0,2)node[above]{$z$};\draw[axis](1,.6)--(.4,1.1)node[above]{$\vct n_{\rm up}$};
\end{tikzpicture}
\caption{向上法向并不是所围封闭体在下曲面上的外法向。}\label{fig:c4-h4-q}
\end{figure}
\begin{solution}
曲面的“向上”法向与闭合体的“向外”法向必须分别确定。图~\ref{fig:c4-h4-q}中抛物面是封闭体的下边界，因此其向外法向朝下，而题目要求的法向朝上；补顶盖使用散度定理时，这个符号差异决定最终相减的顺序。


向上面积元为$(-2kx,-2ky,1)\dd x\dd y$，点积只留下$z=k r^2$，故
\[\Phi_{\rm up}=\iint_Dkr^2\dd A
=2\pi k\int_0^Rr^3\dd r
=2\pi k\left[\frac{r^4}{4}\right]_0^R
=\boxed{\pi kR^4/2}.\]
图~\ref{fig:c4-h4-s}用顶部平面$z=kR^2$封口，所围体积的坐标范围为
\[kr^2\le z\le kR^2,\qquad0\le r\le R,\quad0\le\varphi\le2\pi.\]
在固定$r$处对$z$积分，得到高度$k(R^2-r^2)$，故
\[V=2\pi k\int_0^R(R^2r-r^3)\dd r
=2\pi k\left[\frac{R^2r^2}{2}-\frac{r^4}{4}\right]_0^R
=\frac{\pi kR^4}{2}.\]
散度$\nabla\cdot\vct A=\partial_z z=1$，顶盖向上通量为$kR^2\pi R^2=\pi kR^4$。下曲面的外法向向下，所以
\[\Phi_{\rm top}-\Phi_{\rm up}=V=\pi kR^4/2,
\]
与直接积分一致。闭曲面定理里的外法向不能凭“向上”一词决定，必须相对所围体积判断。


封闭体在$r=R$处高度为零，不再附加一块有限面积的侧壁。顶盖与下曲面由此构成完整闭面。

\begin{figure}[H]\centering
\begin{tikzpicture}[>=Stealth,line cap=round,line join=round]
\fill[gray!12](-1.5,1.35)--(1.5,1.35)--plot[domain=1.5:-1.5]({\x},{.6*(\x)^2})--cycle;\draw[curve,domain=-1.5:1.5]plot(\x,{.6*(\x)^2});\draw[curve](-1.5,1.35)--(1.5,1.35);\draw[axis](0,1.35)--(0,2);\draw[axis](1,.6)--(1.6,.1);\node[right]at(1.6,.1){下曲面外法向};
\end{tikzpicture}
\caption{顶盖外法向向上，下曲面外法向向下，二者围成封闭体。}\label{fig:c4-h4-s}
\end{figure}
\end{solution}

\Needspace{8\baselineskip}
\begin{exercise}[同边界曲面的Stokes积分]\label{ex:c4-h5}
向量场$\vct A=(-y/2,x/2,0)$，边界为圆$x^2+y^2=R^2,z=h$，从上方看逆时针。求环量，并讨论用圆盘或上半球作跨接曲面是否改变结果。 见图~\ref{fig:c4-h5-q}。
\end{exercise}

\begin{figure}[H]\centering
\begin{tikzpicture}[>=Stealth,line cap=round,line join=round]
\draw[curve](0,0)ellipse(1.5 and .6);\draw[axis](1.5,-.1)--(1.5,.2);\draw[axis](0,0)--(0,1.8)node[above]{$+z$};
\end{tikzpicture}
\caption{边界正向从上方看逆时针，圆盘法向向上。}\label{fig:c4-h5-q}
\end{figure}
\begin{solution}
图~\ref{fig:c4-h5-q}的圆周从上方看逆时针，按右手规则应配向上的圆盘法向。
先直接参数化边界：
$\vct r(t)=(R\cos t,R\sin t,h)$，$0\le t\le2\pi$。于是
\[
 \begin{aligned}
 \vct A\cdot\dd\vct r
 &=(-R\sin t/2)(-R\sin t\dd t)
 +(R\cos t/2)(R\cos t\dd t)\\
 &=\frac{R^2}{2}\dd t,
 \end{aligned}
\]
所以
\[\oint_C\vct A\cdot\dd\vct r=\frac{R^2}{2}\int_0^{2\pi}\dd t
=\frac{R^2}{2}[t]_0^{2\pi}=\boxed{\pi R^2}.\]
旋度只有$z$分量非零，等于$\partial_x(x/2)-\partial_y(-y/2)=1/2+1/2=1$，即$\nabla\times\vct A=\vct e_z$；Stokes定理在水平圆盘上同样给出
$\iint \vct e_z\cdot\vct e_z\dd S=\pi R^2$。结果与高度$h$无关。
\begin{figure}[H]\centering
\begin{tikzpicture}[>=Stealth,line cap=round,line join=round]
\draw[curve](-1.6,0)arc(180:0:1.6);\draw[guide](1.6,0)arc(0:180:1.6 and .45);\draw[curve](-1.6,0)arc(180:360:1.6 and .45);
\draw[axis](0,1.6)--(0,2.25)node[above]{$\vct n$};
\draw[axis](1.6,-.12)--(1.6,.2)node[right]{$C$};
\draw[axis](.9,1.3229)--(1.25,1.8374)node[above right]{$\vct n(\theta)$};
\draw[guide](0,0)--(.9,1.3229);\draw(0,.5)arc(90:55.8:.5);\node at(.23,.59){$\theta$};
\node[below]at(0,-.85){球心$(0,0,h)$，球面半径$R$};
\end{tikzpicture}
\caption{上半球的外法向随位置变化，与边界逆时针方向配对；背面半圆用虚线，前面半圆用实线。}\label{fig:c4-h5-s}
\end{figure}
图~\ref{fig:c4-h5-s}中的上半球以$(0,0,h)$为球心，取外法向。局部球坐标下
$n_z=\cos\theta$、$\dd S=R^2\sin\theta\dd\theta\dd\varphi$，
$0\le\theta\le\pi/2$，故
\[
 \begin{aligned}
 \iint_{S_{\rm hemi}}(\nabla\times\vct A)\cdot\vct n\dd S
 &=R^2\int_0^{2\pi}\dd\varphi
 \int_0^{\pi/2}\cos\theta\sin\theta\dd\theta\\
 &=2\pi R^2\left[\frac{\sin^2\theta}{2}\right]_0^{\pi/2}
 =\pi R^2.
 \end{aligned}
\]
这里积分的是有向投影面积，不是半球面积。上半球外法向诱导的边界方向仍为题设的逆时针。
若把半球与圆盘合成闭面，底部圆盘须改取向下法向；由
$\nabla\cdot(\nabla\times\vct A)=0$，闭面通量为零，便得到两张同向跨接曲面的通量相等。
本题场在全空间光滑；若换面穿过场的奇点，须另行处理，不能直接沿用此结论。
\end{solution}

\Needspace{8\baselineskip}
\begin{exercise}[Beta积分求幂型分布的质心]\label{ex:c4-h6}
长度$L$的细杆线密度$\lambda(x)=\lambda_0(x/L)^{p-1}(1-x/L)^{q-1}$，$0<x<L,p,q>0$。求质量与质心，并讨论端点密度发散时质量是否有限。 见图~\ref{fig:c4-h6-q}。
\end{exercise}

\begin{figure}[H]\centering
\begin{tikzpicture}[>=Stealth,line cap=round,line join=round]
\draw[thick](0,0)--(4,0);\draw[axis](0,-.4)--(4.5,-.4)node[right]{$x$};\node[below]at(0,-.4){$0$};\node[below]at(4,-.4){$L$};\draw[curve,domain=.04:3.96]plot(\x,{.3/sqrt((\x/4)*(1-\x/4))});
\end{tikzpicture}
\caption{图取$p=q=1/2$，端点密度发散而总质量有限。}\label{fig:c4-h6-q}
\end{figure}
\begin{solution}
图~\ref{fig:c4-h6-q}的端点密度允许发散，只要幂指数大于$-1$，总质量仍有限。先用$u=x/L$把有限杆长分离出去，余下积分正好具有Beta函数形式；质心则在分子里多乘一个$x$。


令$u=x/L$，则$x=Lu$、$\dd x=L\dd u$，杆的两端变为$u=0,1$。质量与一阶矩分别为
\[M=\lambda_0L\int_0^1u^{p-1}(1-u)^{q-1}\dd u
=\lambda_0LB(p,q),\]
\[\int x\dd m=\lambda_0L^2\int_0^1u^p(1-u)^{q-1}\dd u
=\lambda_0L^2B(p+1,q).\]
使用Beta与Gamma的关系，并分别应用$\Gamma(p+1)=p\Gamma(p)$、$\Gamma(p+q+1)=(p+q)\Gamma(p+q)$，
\[\begin{aligned}
\frac{B(p+1,q)}{B(p,q)}
&=\frac{\Gamma(p+1)\Gamma(q)}{\Gamma(p+q+1)}
\frac{\Gamma(p+q)}{\Gamma(p)\Gamma(q)}\\
&=\frac{p\Gamma(p)}{\Gamma(p)}
\frac{\Gamma(p+q)}{(p+q)\Gamma(p+q)}=\frac p{p+q}.
\end{aligned}\]
故$\boxed{x_C=L p/(p+q)}$。

也可不借助Gamma函数证明这一递推。先在被积函数中插入$u+(1-u)=1$，得
\[B(p+1,q)+B(p,q+1)=B(p,q).\]
再对$u^p(1-u)^q$求导并积分，端点因$p,q>0$均为零，于是
\[0=[u^p(1-u)^q]_0^1=pB(p,q+1)-qB(p+1,q).\]
把第一式的$B(p,q+1)=B(p,q)-B(p+1,q)$代入，便得
$(p+q)B(p+1,q)=pB(p,q)$，与上面的比值相同。

在左端附近，$(1-x/L)^{q-1}$有界且趋于$1$，收敛性由$x^{p-1}$决定：
\[\int_0^\varepsilon x^{p-1}\dd x
=\lim_{\delta\to0^+}\left[\frac{x^p}{p}\right]_\delta^\varepsilon
=\frac{\varepsilon^p}{p}\quad(p>0).\]
所以$0<p<1$虽使密度无界，局部质量仍有限。右端令$v=L-x$，由$q>0$给出相同结论；若某端指数降到$p=0$或$q=0$，便出现对数发散，负指数则出现更强的幂发散。

当$p,q>1$时，密度在两端为零，内部峰可用对数导数确定：
\[\frac{\dd}{\dd u}\ln\lambda
=\frac{p-1}{u}-\frac{q-1}{1-u}=0.\]
通分得$(p-1)(1-u)=(q-1)u$，即
$u_{\rm peak}=(p-1)/(p+q-2)$。二阶导数为
$-(p-1)/u^2-(q-1)/(1-u)^2<0$，故该驻点确为密度峰。
图~\ref{fig:c4-h6-s}取$p=3,q=2$，峰位为$2L/3$，质心却为$3L/5$。一般$p>q$时质心位于中点右侧；$p=q$时反射对称，质心就在中点。密度峰由局部最高值确定，质心由整根杆的一阶矩确定。

\begin{figure}[H]\centering
\begin{tikzpicture}[>=Stealth,line cap=round,line join=round]
\begin{scope}[x=4.5cm,y=12cm]
\draw[axis](0,0)--(1.15,0)node[right]{$x/L$};\draw[axis](0,0)--(0,.21)node[above]{$\lambda/\lambda_0$};
\draw[curve,domain=0:1,samples=101]plot(\x,{(\x)^2*(1-\x)});
\draw[guide](.6,0)--(.6,.144);\fill(.6,.144)circle(1.4pt);
\draw[dashed]({2/3},0)--({2/3},{4/27});\fill({2/3},{4/27})circle(1.4pt);
\node[below left]at(.6,0){$3/5$};\node[below right]at({2/3},0){$2/3$};
\end{scope}
\node[align=left,right]at(5.7,1.6){$p=3,q=2$\\[4pt]质心：$3L/5$\\[4pt]密度峰：$2L/3$};
\end{tikzpicture}
\caption{分别标出一阶矩确定的质心与密度函数的最高点；两条竖线对应不同含义。}\label{fig:c4-h6-s}
\end{figure}
\end{solution}

\Needspace{8\baselineskip}
\begin{exercise}[高斯分布的归一化与二阶矩]\label{ex:c4-h7}
一维密度为$p(x)=A\ee^{-x^2/(2\sigma^2)}$，$x\in\mathbb R,\sigma>0$。求$A$与$\langle x^2\rangle$，并推广到三维独立同分布向量的$\langle r^2\rangle$。 见图~\ref{fig:c4-h7-q}。
\end{exercise}

\begin{figure}[H]\centering
\begin{tikzpicture}[>=Stealth,line cap=round,line join=round]
\begin{scope}[x=0.80000cm,y=3.00000cm]\draw[axis](-3,0)--(3.25,0)node[right]{$x/\sigma$};\draw[axis](0,0)--(0,0.6)node[above]{$\sigma p$};\draw[curve,domain=-3:3,samples=101]plot(\x,{exp(-(\x)^2/2)/sqrt(2*pi)});\end{scope}
\end{tikzpicture}
\caption{归一化取全实轴，不能使用半轴积分而漏掉因子二。}\label{fig:c4-h7-q}
\end{figure}
\begin{solution}
图~\ref{fig:c4-h7-q}的一维密度为偶函数，故$\langle x\rangle=0$。令$a=1/(2\sigma^2)$，记
$Z(a)=\int_{-\infty}^{\infty}\ee^{-ax^2}\dd x=\sqrt{\pi/a}$。
由$AZ(a)=1$，得到
\[
 \boxed{A=\frac1{\sqrt{2\pi}\sigma}}.
\]
在$a>0$处对参数求负导数；高斯衰减保证所需的积分换序成立：
\[\begin{aligned}
-Z'(a)&=\int_{-\infty}^{\infty}x^2\ee^{-ax^2}\dd x
=-\sqrt\pi\,\frac{\dd}{\dd a}a^{-1/2}\\
&=-\sqrt\pi\left(-\frac12a^{-3/2}\right)
=\frac{\sqrt\pi}{2a^{3/2}}.
\end{aligned}\]
因此$\boxed{\langle x^2\rangle=-Z'(a)/Z(a)=1/(2a)=\sigma^2}$。
用本讲的半轴积分记号，该二阶矩也写作$2AG_2(a)$。
\begin{figure}[H]\centering
\begin{tikzpicture}[>=Stealth,line cap=round,line join=round]
\begin{scope}[x=1.00000cm,y=2.25000cm]\draw[axis](0,0)--(4.25,0)node[right]{$r/\sigma$};\draw[axis](0,0)--(0,0.8)node[above]{$\sigma p_r$};\draw[curve,domain=0:4,samples=101]plot(\x,{sqrt(2/pi)*(\x)^2*exp(-(\x)^2/2)});\draw[guide](1.414214,0)--(1.414214,.587051);\end{scope}
\end{tikzpicture}
\caption{三维高斯的径向概率密度包含球壳的$r^2$因子。}\label{fig:c4-h7-s}
\end{figure}
三维独立同分布时，联合密度为
\[
 p_3(\vct r)=(2\pi\sigma^2)^{-3/2}\ee^{-r^2/(2\sigma^2)}.
\]
利用期望的可加性及各坐标的二阶矩，
\[
 \boxed{\langle r^2\rangle
 =\langle x^2\rangle+\langle y^2\rangle+\langle z^2\rangle
 =3\sigma^2}.
\]
图~\ref{fig:c4-h7-s}画的是半径落在窄球壳内的概率密度
$p_r(r)=4\pi r^2p_3(r)$，不是空间点密度$p_3$。
在$r>0$处，
\[
 \frac{\dd}{\dd r}\ln p_r=\frac2r-\frac r{\sigma^2}=0
 \quad\Longrightarrow\quad r_{\rm peak}=\sqrt2\sigma.
\]
它在$r=0$为零、无穷远也趋零，故该驻点为最大值。空间点密度的峰在原点；最常见半径为$\sqrt2\sigma$，均方根半径则为$\sqrt{\langle r^2\rangle}=\sqrt3\sigma$，三者对应不同的量。
\end{solution}

\Needspace{8\baselineskip}
\begin{exercise}[矩形脉冲的连续傅里叶谱]\label{ex:c4-h8}
函数$f(t)=1$在$|t|<\tau/2$内成立，其余为零。按本讲约定求傅里叶变换，确定零点，并说明时间平移与脉宽改变的影响。 见图~\ref{fig:c4-h8-q}。
\end{exercise}

\begin{figure}[H]\centering
\begin{tikzpicture}[>=Stealth,line cap=round,line join=round]
\draw[axis](-2.4,0)--(2.5,0)node[right]{$t$};\draw[axis](0,-.2)--(0,1.7)node[above]{$f$};\draw[curve](-2,0)--(-1,0)--(-1,1)--(1,1)--(1,0)--(2,0);\node[below]at(-1,0){$-\tau/2$};\node[below]at(1,0){$\tau/2$};
\end{tikzpicture}
\caption{对称矩形脉冲的变换为实偶函数。}\label{fig:c4-h8-q}
\end{figure}
\begin{solution}
矩形脉冲可由有限区间直接积分。图~\ref{fig:c4-h8-q}关于$t=0$对称，所以变换的虚部抵消，频谱为实偶函数；中心移位才会产生线性相位，不能把它误算成幅度变化。


对$\omega\ne0$，定义积分为
\[\begin{aligned}
F(\omega)&=\int_{-\tau/2}^{\tau/2}\ee^{-\ii\omega t}\dd t
=\left[\frac{\ee^{-\ii\omega t}}{-\ii\omega}\right]_{-\tau/2}^{\tau/2}\\
&=\frac{\ee^{-\ii\omega\tau/2}-\ee^{\ii\omega\tau/2}}{-\ii\omega}
=\frac{2\sin(\omega\tau/2)}{\omega}.
\end{aligned}\]
在$\omega=0$处，原积分直接等于$\tau$，也等于商式的连续极限。定义$\operatorname{sinc}u=\sin u/u$并连续定义$\operatorname{sinc}0=1$，于是
\[\boxed{F(\omega)=\tau\operatorname{sinc}(\omega\tau/2)}.
\]
$F(0)=\tau$等于时间面积，并非谱零点。非零谱零点在$\omega=2\pi n/\tau$，$n\in\mathbb Z\setminus\{0\}$。图~\ref{fig:c4-h8-s}的主瓣位于首对零点$\pm2\pi/\tau$之间，宽度为$4\pi/\tau$，因此脉冲越短，谱越宽。

将脉冲移到中心$t_0$时，令$u=t-t_0$，积分限仍为$\pm\tau/2$，且$\dd t=\dd u$。由指数分解，
\[\int_{t_0-\tau/2}^{t_0+\tau/2}\ee^{-\ii\omega t}\dd t
=\ee^{-\ii\omega t_0}\int_{-\tau/2}^{\tau/2}\ee^{-\ii\omega u}\dd u
=\ee^{-\ii\omega t_0}F(\omega).\]
所以模长不变，相位增加$-\omega t_0$。孤立跳跃点的赋值不影响变换。

\begin{figure}[H]\centering
\begin{tikzpicture}[>=Stealth,line cap=round,line join=round]
\begin{scope}[x=0.55000cm,y=1.05882cm]\draw[axis](-6.3,0)--(6.55,0)node[right]{$\omega\tau/2$};\draw[axis](0,-0.4)--(0,1.3)node[above]{$F/\tau$};\draw[curve,domain=-6.3:6.3,samples=101]plot(\x,{sin(\x r)/\x});\end{scope}
\end{tikzpicture}
\caption{谱零点等间隔，主瓣宽度与时间脉宽成反比。}\label{fig:c4-h8-s}
\end{figure}
\end{solution}

\Needspace{8\baselineskip}
\begin{exercise}[指数序列的周期频谱]\label{ex:c4-h9}
序列$x_n=a^n$（$n\ge0$），$n<0$时为零，$|a|<1$。求DTFT，并用逆变换说明为何只积一个周期。 见图~\ref{fig:c4-h9-q}。
\end{exercise}

\begin{figure}[H]\centering
\begin{tikzpicture}[>=Stealth,line cap=round,line join=round]
\draw[axis](-.5,0)--(4.5,0)node[right]{$n$};\draw[axis](0,-.2)--(0,1.5)node[above]{$x_n$};\foreach\n in{0,...,4}{\draw[curve](\n,0)--(\n,{.65^\n});\fill(\n,{.65^\n})circle(1.5pt);}
\end{tikzpicture}
\caption{示意$a=.65$的因果指数序列。}\label{fig:c4-h9-q}
\end{figure}
\begin{solution}
离散时间的指标$n$为整数，因此频率增加$2\pi$后每一项相位都不变。图~\ref{fig:c4-h9-q}表示因果衰减序列；$|a|<1$保证绝对收敛，随后才可用无穷等比级数求和。


令$z=a\ee^{-\ii\omega}$。有限和满足$(1-z)\sum_{n=0}^Nz^n=1-z^{N+1}$。因$|z|=|a|<1$，当$N\to\infty$时$z^{N+1}\to0$，故
\[\boxed{X(\omega)=\sum_{n=0}^\infty(a\ee^{-\mathrm i\omega})^n=\frac1{1-a\ee^{-\mathrm i\omega}}}.
\]
因为$\ee^{-\ii(\omega+2\pi)n}=\ee^{-\ii\omega n}$，所以$X(\omega+2\pi)=X(\omega)$。


把逆变换中的求和指标记为$m$，可避免与待恢复的$n$混淆：
\[x_n=\sum_{m=0}^\infty a^m\frac1{2\pi}
\int_{-\pi}^{\pi}\ee^{\mathrm i(n-m)\omega}\dd\omega.\]
当$n=m$时，归一化积分为$[\omega]_{-\pi}^\pi/(2\pi)=1$；当整数$\ell=n-m\ne0$时，
\[\frac1{2\pi}\int_{-\pi}^{\pi}\ee^{\ii\ell\omega}\dd\omega
=\frac1{2\pi}\left[\frac{\ee^{\ii\ell\omega}}{\ii\ell}\right]_{-\pi}^\pi
=\frac{\sin(\ell\pi)}{\pi\ell}=0.\]
因此$n\ge0$时只留下$a^n$，$n<0$时没有可留下的$m$，恢复$x_n=0$。这里求和与积分的交换由$\sum|a|^m<\infty$保证。图~\ref{fig:c4-h9-s}中相邻谱峰间距为$2\pi$，选择任意一段完整周期均可反演；若在整个实轴重复积分，每一周期会给出相同贡献。

一般复数$a$满足$|X|=[1+|a|^2-2\Re(a\ee^{-\ii\omega})]^{-1/2}$。若$a=\rho\ee^{\ii\phi}$，则
\[|X(\omega)|=[1+\rho^2-2\rho\cos(\omega-\phi)]^{-1/2}.\]
在$0<\rho<1$时，幅谱的峰值位于$\omega=\phi\pmod{2\pi}$，峰高$1/(1-\rho)$；$a=0$时频谱恒为$1$，没有独立的峰位。

图~\ref{fig:c4-h9-s}取实正$a=0.65$，此时模谱为$[1+a^2-2a\cos\omega]^{-1/2}$，峰在零频率，且因序列实值而有$X(-\omega)=X(\omega)^*$。一般非实$a$仍保持周期性，却没有这一共轭对称；这两个性质来自不同条件。

\begin{figure}[H]\centering
\begin{tikzpicture}[>=Stealth,line cap=round,line join=round]
\begin{scope}[x=.7cm,y=.9cm]
\draw[axis](-6.6,0)--(6.65,0)node[right]{$\omega$};\draw[axis](0,0)--(0,3.25)node[above]{$|X|$};
\draw[curve,domain=-6.283185:6.283185,samples=201]plot(\x,{1/sqrt(1+.65^2-1.3*cos(\x r))});
\foreach\x/\lab in{-6.283185/-2\pi,-3.141593/-\pi,3.141593/\pi,6.283185/2\pi}{\draw(\x,0)--(\x,-.08)node[below]{$\lab$};}
\draw[<->](0,3.9)--({2*pi},3.9)node[midway,above]{$2\pi$};\end{scope}
\end{tikzpicture}
\caption{实参数$a=0.65$的DTFT模谱。标出相邻峰之间的周期；一般复参数会改变峰的横向位置。}\label{fig:c4-h9-s}
\end{figure}
\end{solution}

\Needspace{8\baselineskip}
\begin{exercise}[DFT的频率格点与混叠]\label{ex:c4-h10}
长度$N$的序列$x_n=\cos(2\pi mn/N)$，其中$0<m<N/2$且$m$为整数。求其DFT；若以采样间隔$\Delta t$采样连续正弦，说明频率为何只能在相差$2\pi/\Delta t$的等价类中辨认。 见图~\ref{fig:c4-h10-q}。
\end{exercise}

\begin{figure}[H]\centering
\begin{tikzpicture}[>=Stealth,line cap=round,line join=round]
\draw[axis](-.3,0)--(5,0)node[right]{$n$};\draw[axis](0,-1.4)--(0,1.5)node[above]{$x_n$};\foreach\n in{0,...,7}{\draw[curve]({.6*\n},0)--({.6*\n},{cos(90*\n)});\fill({.6*\n},{cos(90*\n)})circle(1.5pt);}
\end{tikzpicture}
\caption{示意$N=8,m=2$，采样只保留格点上的值。}\label{fig:c4-h10-q}
\end{figure}
\begin{solution}
图~\ref{fig:c4-h10-q}的余弦在窗口内完成整数个周期。将它分为两个复指数，代入DFT，
\[X_k=\frac12\sum_{n=0}^{N-1}\ee^{2\pi\ii(m-k)n/N}
+\frac12\sum_{n=0}^{N-1}\ee^{-2\pi\ii(m+k)n/N}.\]
有限正交和可由几何级数求出：$z=\ee^{2\pi\ii\ell/N}\ne1$时，$\sum_{n=0}^{N-1}z^n=(1-z^N)/(1-z)=0$；$z=1$时每项为一，总和为$N$。因此
\[\boxed{X_m=X_{N-m}=N/2,\qquad\text{其余}X_k=0}.\]
图~\ref{fig:c4-h10-s}中的两点对应正负频率；所取$N=8,m=2$给出$k=2,6$，各点高度为$4$，时域样点为$1,0,-1,0,1,0,-1,0$。条件$0<m<N/2$保证两点不同。直流$m=0$及偶数$N$的Nyquist点$m=N/2$会使两个指数重合，必须先合并，不能仍分成两个互异的$N/2$谱峰。

对物理频率$\Omega$，采样只保留$t_n=n\Delta t$。由
\[\ee^{\ii(\Omega+2\pi j/\Delta t)n\Delta t}=\ee^{\ii\Omega n\Delta t}
\]
可知，连续复指数的频率相差整数倍采样角频率时不可辨。实余弦还具有反号等价：
\[\cos(\Omega t+\phi)=\cos(-\Omega t-\phi),\qquad
(\Omega,\phi)\sim(-\Omega+2\pi j/\Delta t,-\phi)
\]
（后式指采样序列相同）。对正幅度的实单频信号，可选非负代表区间$0\le\Omega\le\pi/\Delta t$；开区间内的一般信号保留幅度与相位信息，两个端点则有特殊退化，例如Nyquist点只看到$(-1)^n\cos\phi$，不能独立恢复任意相位。实际消除混叠还需在采样前限制输入频带。

有限窗口造成的泄漏是另一问题。不在DFT格点上的频率乘上矩形观测窗后，会向多个格点分散；整周期入窗只消除了这一例中的泄漏，没有解除采样本身的频率等价。仅补零会细化频谱显示点，不能增加原始观测时长或消除混叠。

\begin{figure}[H]\centering
\begin{tikzpicture}[>=Stealth,line cap=round,line join=round]
\draw[axis](-.3,0)--(5.2,0)node[right]{$k$};\draw[axis](0,-.2)--(0,2)node[above]{$|X_k|$};\draw[curve](1.2,0)--(1.2,1.5) (3.6,0)--(3.6,1.5);\fill(1.2,1.5)circle(2pt);\fill(3.6,1.5)circle(2pt);\node[below]at(1.2,0){$m$};\node[below]at(3.6,0){$N-m$};\node[above]at(1.2,1.5){$N/2$};
\end{tikzpicture}
\caption{整周期入窗时只出现一对互为共轭的离散谱点。}\label{fig:c4-h10-s}
\end{figure}
\end{solution}

\Needspace{8\baselineskip}
\begin{exercise}[延迟阶跃的拉普拉斯响应]\label{ex:c4-h11}
一阶系统$\dot y+ay=F_0H(t-\tau)$，$y(0)=0,a>0,\tau>0$。求拉普拉斯变换及全时间响应，$H$为单位阶跃。 见图~\ref{fig:c4-h11-q}。
\end{exercise}

\begin{figure}[H]\centering
\begin{tikzpicture}[>=Stealth,line cap=round,line join=round]
\draw[axis](-.3,0)--(4.5,0)node[right]{$t$};\draw[axis](0,-.2)--(0,1.8);\draw[curve](0,0)--(1.5,0)--(1.5,1.2)--(4,1.2);\node[below]at(1.5,0){$\tau$};\node[above]at(3,1.2){$F_0$};
\end{tikzpicture}
\caption{外源在延迟时间后开启，之前系统保持静止。}\label{fig:c4-h11-q}
\end{figure}
\begin{solution}
图~\ref{fig:c4-h11-q}中的驱动在$t=\tau$开启，其变换为
\[\int_0^\infty H(t-\tau)\ee^{-st}\dd t
=\int_\tau^\infty\ee^{-st}\dd t
=\left[-\frac{\ee^{-st}}s\right]_\tau^\infty
=\frac{\ee^{-s\tau}}s\quad(\Re s>0).\]
对方程作单边拉普拉斯变换，得
\[
 (sY-y(0))+aY=\frac{F_0\ee^{-s\tau}}s,\qquad
 \boxed{Y(s)=\frac{F_0\ee^{-s\tau}}{s(s+a)}}.
\]
先将未延迟部分分解为
\[
 \frac{F_0}{s(s+a)}
 =\frac{F_0}{a}\left(\frac1s-\frac1{s+a}\right).
\]
其反演为$(F_0/a)(1-\ee^{-at})$。再用$\ee^{-s\tau}$对应的因果时移，
\[
 \boxed{y(t)=\frac{F_0}{a}
 [1-\ee^{-a(t-\tau)}]H(t-\tau)}.
\]
即$t<\tau$时$y=0$，$t>\tau$时取括号内的指数响应；开启前若处于题设静止状态，负时间也为零。阶跃因子不能省略，否则公式会在驱动开启前给出非零响应。
\begin{figure}[H]\centering
\begin{tikzpicture}[>=Stealth,line cap=round,line join=round]
\begin{scope}[x=1.60000cm,y=1.28571cm]\draw[axis](0,0)--(1.75,0);\draw[axis](0,0)--(0,1.4)node[above]{$ay/F_0$};\draw[curve,domain=0:1.5,samples=101]plot(\x,{0});\draw[curve,domain=1.5:5,samples=101]plot(\x,{1-exp(-(\x-1.5))});\draw[axis](1.5,0)--(5.2,0)node[right]{$t$};\draw[guide](1.5,0)--(1.5,1.2);\node[below]at(1.5,0){$\tau$};\end{scope}
\end{tikzpicture}
\caption{开启前响应为零，开启后按同一时间常数趋于静态值。}\label{fig:c4-h11-s}
\end{figure}
在$t>\tau$处求导，得$\dot y=F_0\ee^{-a(t-\tau)}$，故
$\dot y+ay=F_0$，且$t\to\infty$时$y\to F_0/a$。
为确定启动连接条件，在$(\tau-\epsilon,\tau+\epsilon)$上积分原方程：
\[
 y(\tau+\epsilon)-y(\tau-\epsilon)
 +a\int_{\tau-\epsilon}^{\tau+\epsilon}y\dd t
 =F_0\int_{\tau-\epsilon}^{\tau+\epsilon}H(t-\tau)\dd t.
\]
后两个积分随$\epsilon\to0$趋零，所以$y$连续；结合开启前$y=0$，得
$y(\tau^-)=y(\tau^+)=0$。
由两侧方程进一步有$\dot y(\tau^-)=0$、$\dot y(\tau^+)=F_0$，
与图~\ref{fig:c4-h11-s}中连续起步、斜率突变的曲线一致。
$H(0)$的取值不改变这些单侧极限。若将外源换成冲量，右端的小区间积分有有限值，状态才会发生跃变。
\end{solution}

\Needspace{8\baselineskip}
\begin{exercise}[从卷积到串联系统]\label{ex:c4-h12}
两个因果一阶系统的单位冲量响应为$g_1(t)=\ee^{-at}H(t)$、$g_2(t)=\ee^{-bt}H(t)$，$a,b>0$。求串联响应，并讨论$a=b$极限及该极限与重复特征根的关系。 见图~\ref{fig:c4-h12-q}。
\end{exercise}

\begin{figure}[H]\centering
\begin{tikzpicture}[>=Stealth,line cap=round,line join=round]
\node[draw,minimum width=1.5cm,minimum height=.8cm](A)at(0,0){$g_1$};\node[draw,minimum width=1.5cm,minimum height=.8cm](B)at(3,0){$g_2$};\draw[axis](-2,0)--(A)node[midway,above]{$\delta(t)$};\draw[axis](A)--(B);\draw[axis](B)--(5,0)node[right]{$g$};
\end{tikzpicture}
\caption{串联的中间输出再驱动第二个系统，时域运算为卷积。}\label{fig:c4-h12-q}
\end{figure}
\begin{solution}
图~\ref{fig:c4-h12-q}中的两个核均为因果指数。固定观察时刻$t\ge0$，只有$0\le\tau\le t$的输入同时落在两核的非零区间，所以
\[\begin{aligned}
g(t)&=\int_0^t\ee^{-a\tau}\ee^{-b(t-\tau)}\dd\tau\\
&=\ee^{-bt}\int_0^t\ee^{(b-a)\tau}\dd\tau\\
&=\ee^{-bt}\left[\frac{\ee^{(b-a)\tau}}{b-a}\right]_0^t
=\boxed{\frac{\ee^{-at}-\ee^{-bt}}{b-a}}\quad(a\ne b).
\end{aligned}\]
负时间响应为零。若$a=b$，被积函数恒为$\ee^{-at}$，所以正时间内
$g(t)=\ee^{-at}\int_0^t\dd\tau=\ee^{-at}[\tau]_0^t=t\ee^{-at}$，完整因果形式为
\[\boxed{g(t)=t\ee^{-at}H(t)}.\]
两种结果都满足$g(0^+)=0,g'(0^+)=1$，长时趋零。变换域中$G(s)=1/[(s+a)(s+b)]$；两极点合并后成为$1/(s+a)^2$，反演中的时间因子$t$与常系数方程重根的第二解具有相同来源。

对$a\ne b$的响应求导，
\[g'(t)=\frac{-a\ee^{-at}+b\ee^{-bt}}{b-a}.\]
峰点条件为$b\ee^{-bt_*}=a\ee^{-at_*}$。两边均为正，取对数得
$\ln b-bt_*=\ln a-at_*$，即$(b-a)t_*=\ln(b/a)$，于是
\[\boxed{t_* =\frac{\ln(b/a)}{b-a}>0}.\]
分子分母同号，且该方程只有一个正根；结合初始上升与长时衰减，可知它是唯一峰值。$b\to a$时，$\ln(b/a)\sim(b-a)/a$，故$t_*\to1/a$。在重根情形直接用乘积法则也有
$g'(t)=\ee^{-at}-at\ee^{-at}=\ee^{-at}(1-at)$，其零点同为$1/a$。图~\ref{fig:c4-h12-s}画的是$a=b$，标出的$1/a$只适用于这一情形。

闭式交换$a,b$后不变，说明这两个线性时不变系统可交换串联次序。若系统参数随时间改变，响应一般不再只依赖$t-\tau$，上述卷积表达与交换结论便须重新分析。

\begin{figure}[H]\centering
\begin{tikzpicture}[>=Stealth,line cap=round,line join=round]
\begin{scope}[x=1.00000cm,y=3.27273cm]\draw[axis](0,0)--(5.25,0)node[right]{$at$};\draw[axis](0,0)--(0,0.55)node[above]{$ag$};\draw[curve,domain=0:5,samples=101]plot(\x,{\x*exp(-\x)});\fill(1,.367879)circle(1.5pt)node[above]{$t=1/a$};\end{scope}
\end{tikzpicture}
\caption{两个相同衰减率串联，产生带时间因子的响应。}\label{fig:c4-h12-s}
\end{figure}
\end{solution}

\Needspace{8\baselineskip}
\begin{exercise}[方波系数的均方和]\label{ex:c4-h13}
用本讲方波的傅里叶系数和Parseval关系求$\sum_{j=0}^\infty(2j+1)^{-2}$，再求全部正整数平方倒数和。 见图~\ref{fig:c4-h13-q}。
\end{exercise}

\begin{figure}[H]\centering
\begin{tikzpicture}[>=Stealth,line cap=round,line join=round]
\draw[axis](0,0)--(5,0)node[right]{$n$};\draw[axis](0,0)--(0,2)node[above]{$b_n$};\foreach\n in{1,3,5,7}{\draw[curve]({.6*\n},0)--({.6*\n},{1.5/\n});\fill({.6*\n},{1.5/\n})circle(1.5pt);}
\end{tikzpicture}
\caption{方波只有奇次正弦系数，均方值为一。}\label{fig:c4-h13-q}
\end{figure}
\begin{solution}
方波除跳点外满足$f^2=1$，其周期均方值为
\[\frac1{2\pi}\int_{-\pi}^{\pi}f^2\dd t
=\frac1{2\pi}[t]_{-\pi}^\pi=1.\]
跳点的赋值不影响积分。图~\ref{fig:c4-h13-q}的非零系数只有$b_{2j+1}=4/[\pi(2j+1)]$。由正弦函数的正交关系，每一项的均方贡献为$b_{2j+1}^2/2$，因此Parseval等式给出
\[1=\frac12\sum_{j=0}^\infty\frac{16}{\pi^2(2j+1)^2},\qquad
\boxed{\sum_{j=0}^\infty\frac1{(2j+1)^2}=\frac{\pi^2}{8}}.\]
若记全部平方倒数和为$S$，偶数项之和为
\[\sum_{j=1}^\infty\frac1{(2j)^2}=\frac14S.\]
余下奇数项为$3S/4=\pi^2/8$，所以$\boxed{S=\pi^2/6}$。

截断到第$2N+1$次谐波时，已保留的均方贡献是图~\ref{fig:c4-h13-s}中的累计值。误差与部分和正交，误差均方等于总均方减去已保留部分：
\[\frac1{2\pi}\int_{-\pi}^{\pi}|f-S_{2N+1}|^2\dd t
=1-\frac8{\pi^2}\sum_{j=0}^N\frac1{(2j+1)^2}.\]
右侧随$N$增大趋于零。跳点附近仍可能有局部过冲，但误差较大的区域逐渐收窄，因而不妨碍均方收敛。

\begin{figure}[H]\centering
\begin{tikzpicture}[>=Stealth,line cap=round,line join=round]
\draw[axis](0,0)--(5,0)node[right]{$N$};\draw[axis](0,0)--(0,2)node[above]{累计均方};\draw[guide](0,1.5)--(4.5,1.5)node[right]{$1$};\draw[curve](0,1.21585)--(1,1.35094)--(2,1.39957)--(3,1.42439)--(4,1.4394);\node at(2.8,.6){$\frac8{\pi^2}\sum_{j=0}^N(2j+1)^{-2}$};
\end{tikzpicture}
\caption{累计均方贡献单调趋于方波的总均方值。}\label{fig:c4-h13-s}
\end{figure}
\end{solution}

\Needspace{8\baselineskip}
\begin{exercise}[高斯函数的傅里叶变换]\label{ex:c4-h14}
取$f(t)=\ee^{-at^2}$，$a>0$。不使用复平面积分，由对频率求导和分部积分求$F(\omega)$，说明时域宽度与频域宽度的关系。 见图~\ref{fig:c4-h14-q}。
\end{exercise}

\begin{figure}[H]\centering
\begin{tikzpicture}[>=Stealth,line cap=round,line join=round]
\begin{scope}[x=0.70000cm,y=1.38462cm]\draw[axis](-3,0)--(3.25,0)node[right]{$t$};\draw[axis](0,0)--(0,1.3)node[above]{$f$};\draw[curve,domain=-3:3,samples=101]plot(\x,{exp(-.5*(\x)^2)});\end{scope}
\end{tikzpicture}
\caption{高斯函数在无穷远充分衰减，可交换参数求导并作分部积分。}\label{fig:c4-h14-q}
\end{figure}
\begin{solution}
图~\ref{fig:c4-h14-q}中的高斯衰减使$|t|\ee^{-at^2}$可积，因而可在积分号内对$\omega$求导。又$t\ee^{-at^2}=-(2a)^{-1}\partial_t\ee^{-at^2}$，故
\[\begin{aligned}
F'(\omega)&=-\ii\int_{-\infty}^{\infty}t\ee^{-at^2}\ee^{-\ii\omega t}\dd t\\
&=\frac{\ii}{2a}\int_{-\infty}^{\infty}\partial_t(\ee^{-at^2})\ee^{-\ii\omega t}\dd t\\
&=\frac{\ii}{2a}\left\{
\left[\ee^{-at^2}\ee^{-\ii\omega t}\right]_{-\infty}^{\infty}
+\ii\omega\int_{-\infty}^{\infty}\ee^{-at^2}\ee^{-\ii\omega t}\dd t\right\}\\
&=-\frac\omega{2a}F(\omega).
\end{aligned}\]
边界项为零，最后的负号来自$\ii^2=-1$。这是一阶线性方程，积分因子为
\[\mu(\omega)=\exp\left(\int\frac\omega{2a}\dd\omega\right)
=\exp\left(\frac{\omega^2}{4a}\right).\]
用乘积法则合并，避免对未知复函数取对数：
\[\frac{\dd}{\dd\omega}\left[\ee^{\omega^2/(4a)}F(\omega)\right]=0.
\]
从$0$到$\omega$积分，$\mu(\omega)F(\omega)-\mu(0)F(0)=0$。由$\mu(0)=1$及高斯积分$F(0)=\sqrt{\pi/a}$确定常数，得到
\[\boxed{F(\omega)=\sqrt{\frac\pi a}\,\ee^{-\omega^2/(4a)}}.\]

图~\ref{fig:c4-h14-s}的频谱仍为高斯。若按幅度包络定义宽度，写作
\[f(t)=\ee^{-t^2/(2\sigma_t^2)},\qquad
F(\omega)=F(0)\ee^{-\omega^2/(2\sigma_\omega^2)},\]
比较指数得$\sigma_t=1/\sqrt{2a}$、$\sigma_\omega=\sqrt{2a}$，因此$\sigma_t\sigma_\omega=1$。若将$|f|^2$、$|F|^2$分别归一化为密度，其标准差各再除以$\sqrt2$，二者乘积为$1/2$。宽度定义不同，数值因子也随之改变。

参数求导需要可积控制，分部积分需要端点衰减；本题的高斯因子同时保证两项条件。不能把推导直接移用于普通积分意义下不收敛的振荡函数。

\begin{figure}[H]\centering
\begin{tikzpicture}[>=Stealth,line cap=round,line join=round]
\begin{scope}[x=0.60000cm,y=1.38462cm]\draw[axis](-4,0)--(4.25,0)node[right]{$\omega$};\draw[axis](0,0)--(0,1.3)node[above]{$F/\sqrt{\pi/a}$};\draw[curve,domain=-4:4,samples=101]plot(\x,{exp(-.25*(\x)^2)});\end{scope}
\end{tikzpicture}
\caption{频谱宽度与时域宽度相反变化，归一化峰值由高斯积分确定。}\label{fig:c4-h14-s}
\end{figure}
\end{solution}

\Needspace{8\baselineskip}
\begin{exercise}[终值定理的适用范围]\label{ex:c4-h15}
分别考察$F_1(s)=1/[s(s+1)]$、$F_2(s)=1/(s^2+1)$、$F_3(s)=1/(s-1)$的因果反演。比较$\lim_{s\to0}sF_i(s)$与真实时域终值。 见图~\ref{fig:c4-h15-q}。
\end{exercise}

\begin{figure}[H]\centering
\begin{tikzpicture}[>=Stealth,line cap=round,line join=round]
\draw[axis](-2,0)--(2,0)node[right]{$\Re s$};\draw[axis](0,-1.6)--(0,1.6)node[above]{$\Im s$};\fill(-1,0)circle(1.5pt)node[above]{$-1$};\fill(1,0)circle(1.5pt)node[above]{$1$};\fill(0,1)circle(1.5pt);\fill(0,-1)circle(1.5pt);
\end{tikzpicture}
\caption{左半平面、虚轴与右半平面极点对应不同的长时行为。}\label{fig:c4-h15-q}
\end{figure}
\begin{solution}
对有理拉普拉斯变换，常用的终值定理充分条件为：
$sF(s)$的所有极点均严格位于左半平面。满足这一条件时，才可用
$\lim_{t\to\infty}f(t)=\lim_{s\to0^+}sF(s)$。
图~\ref{fig:c4-h15-q}标出本题需要区别的三类极点。

第一式先作部分分式分解：
\[
 F_1=\frac1s-\frac1{s+1},\qquad
 f_1(t)=(1-\ee^{-t})H(t).
\]
$sF_1=1/(s+1)$只有极点$-1$，满足条件；真实终值和变换域极限均为$1$。

第二式反演为$f_2(t)=\sin t\,H(t)$。虽然
$\lim_{s\to0^+}s/(s^2+1)=0$，但$sF_2$有虚轴极点$\pm\ii$，
振荡不衰减，所以时间终值不存在，不能套用定理。
\begin{figure}[H]\centering
\begin{tikzpicture}[>=Stealth,line cap=round,line join=round]
\begin{scope}[x=1.00000cm,y=1.00000cm]\draw[axis](0,0)--(6.25,0)node[right]{$t$};\draw[axis](0,-1.3)--(0,3)node[above]{$f$};\draw[curve,domain=0:6,samples=101]plot(\x,{sin(\x r)});\draw[dashed,domain=0:1.1]plot(\x,{exp(\x)});\node[right]at(2.3,2.4){实线$\sin t$};\node[right]at(2.3,1.8){虚线$\ee^t$};\end{scope}
\end{tikzpicture}
\caption{持续振荡和指数增长都没有有限终值，代数极限不能直接解释成时域极限。}\label{fig:c4-h15-s}
\end{figure}
第三式反演为$f_3(t)=\ee^tH(t)$，其$sF_3=s/(s-1)$有右半平面极点$1$，不满足条件，真实响应发散。
形式上的代数极限仍为$0$，却不代表时间终值。
这里拉普拉斯积分仅在$\operatorname{Re}s>1$收敛；有理表达式可以延拓到$s=0$附近，但在那里取值已经不是在原收敛积分内取极限。

\end{solution}

\Needspace{8\baselineskip}
\begin{exercise}[从非均匀球体的转动惯量积分到Beta函数]\label{ex:c4-h16}
半径$R$的实心球密度$\rho(r)=\rho_0(1-r^2/R^2)^p$，$p>-1$。求质量与绕直径的转动惯量，以Beta函数表示，并化简$I/(MR^2)$。 见图~\ref{fig:c4-h16-q}。
\end{exercise}

\begin{figure}[H]\centering
\begin{tikzpicture}[>=Stealth,line cap=round,line join=round]
\fill[gray!12](0,0)circle(1.5);\draw[curve](0,0)circle(1.5);\draw[curve](0,0)circle(1);\draw[axis](0,-1.9)--(0,2)node[above]{$z$};\draw[axis](0,0)--(1,0)node[midway,above]{$r$};
\end{tikzpicture}
\caption{按球壳分层，每层密度只依赖径向距离。}\label{fig:c4-h16-q}
\end{figure}
\begin{solution}
球对称密度可先积掉角变量。质量积分的角向因子为$4\pi$；图~\ref{fig:c4-h16-s}中的质量元$P$距转轴为$r\sin\theta$，惯量还要乘$r^2\sin^2\theta$，相应角积分为
\[\int\sin^2\theta\dd\Omega
=2\pi\int_0^\pi(1-\cos^2\theta)\sin\theta\dd\theta
=2\pi\left[-\cos\theta+\frac{\cos^3\theta}{3}\right]_0^\pi
=\frac{8\pi}3.\]
其中一个$\sin\theta$来自球坐标体积元，另外两个来自转动半径的平方。因而
\[M=4\pi\rho_0\int_0^Rr^2(1-r^2/R^2)^p\dd r,\]
\[I=\frac{8\pi\rho_0}3\int_0^Rr^4(1-r^2/R^2)^p\dd r.\]
\begin{figure}[H]\centering
\begin{tikzpicture}[>=Stealth,line cap=round,line join=round]
\draw[curve](0,0)circle(1.6);\draw[axis](0,-1.9)--(0,2.1)node[above]{$z$};
\coordinate(P)at(1.28,.96);\fill(P)circle(1.5pt)node[above right]{$P$};
\draw[axis](0,0)--(P)node[midway,below right]{$r$};\draw[guide](P)--(0,.96);
\draw[<->](0,1.23)--(1.28,1.23)node[midway,above]{$r\sin\theta$};
\draw(0,.48)arc(90:36.87:.48);\node at(.25,.56){$\theta$};
\node[right]at(2.1,.4){$\dd I=(r\sin\theta)^2\dd m$};
\end{tikzpicture}
\caption{半径$r$由球心量到质量元$P$，转动半径则是$P$到$z$轴的垂直距离$r\sin\theta$。}\label{fig:c4-h16-s}
\end{figure}

令$u=r^2/R^2$，则$r=R\sqrt u$、$\dd r=R u^{-1/2}\dd u/2$，所以
\[r^2\dd r=\frac{R^3}2u^{1/2}\dd u,\qquad
r^4\dd r=\frac{R^5}2u^{3/2}\dd u.\]
两式积分域均变为$0\le u\le1$，分别识别为Beta积分：
\[\boxed{M=2\pi\rho_0R^3B(3/2,p+1)},\qquad
I=\frac{4\pi\rho_0R^5}3B(5/2,p+1).\]
端点$u=1$附近的幂为$p$，故收敛条件为题设的$p>-1$。再用递推关系$B(a+1,b)/B(a,b)=a/(a+b)$，
\[\frac I{MR^2}=\frac23\frac{B(5/2,p+1)}{B(3/2,p+1)}
=\frac23\frac{3/2}{p+5/2}
=\boxed{\frac2{2p+5}}.\]
图~\ref{fig:c4-h16-q}中$p$增大时密度向中心集中，惯量比值下降；$p=0$恢复均匀实心球的$2/5$。当$p\to-1^+$时，比值趋于薄球壳的$2/3$。若保持$\rho_0$固定，总质量会发散；趋于薄壳的是质量归一化后的径向分布。若要保持总质量有限且固定，应相应调整密度幅值。


\end{solution}

\begin{exercise}[正负两次推力后的余振]\label{ex:c4-v2}

质量$m$、劲度系数$k$的无阻尼振子从平衡点静止出发。取$F_0>0$，先在$0<t<\tau$施加恒力$F_0$，再在$\tau<t<2\tau$施加$-F_0$，随后撤力。两段外力的总冲量为零。求撤力后的振幅、能量以及完全没有余振的$\tau$，并判断总冲量为零是否足以使振子复原。 见图~\ref{fig:c4-v2-q}。
\end{exercise}

\begin{figure}[H]\centering
\begin{tikzpicture}[>=Stealth,line cap=round,line join=round]

\draw[axis](-.3,0)--(5.7,0)node[right]{$t$};\draw[axis](0,-1.5)--(0,1.6)node[above]{$F$};
\draw[curve](0,0)--(0,1.1)--(1.8,1.1)--(1.8,-1.1)--(3.6,-1.1)--(3.6,0)--(5.3,0);
\node[above]at(.8,1.1){$F_0$};\node[below]at(2.7,-1.1){$-F_0$};\node[below left]at(1.8,0){$\tau$};\node[below right]at(3.6,0){$2\tau$};

\end{tikzpicture}
\caption{两段面积大小相同、符号相反，但振子在两段内所处的相位不同。}\label{fig:c4-v2-q}
\end{figure}
\begin{solution}
图~\ref{fig:c4-v2-q}的三个跃变时刻为$0,\tau,2\tau$。令$\omega=\sqrt{k/m}$，单位阶跃记为$H$，则
\[F(t)=F_0[H(t)-2H(t-\tau)+H(t-2\tau)].\]
对零初态，Laplace变换给出
\[(ms^2+k)X(s)=\frac{F_0}{s}
\bigl(1-2\ee^{-s\tau}+\ee^{-2s\tau}\bigr),\]
将右侧配成平方并用$k=m\omega^2$，有
\[X(s)=\frac{F_0}{m}\frac{(1-\ee^{-s\tau})^2}{s(s^2+\omega^2)}.
\]
单个恒力阶跃在正时间内满足$\ddot x+\omega^2x=F_0/m$。取常数特解$F_0/k$，齐次方程的特征根为$\pm\ii\omega$，故
\[x=\frac{F_0}{k}+A\cos\omega t+B\sin\omega t,\qquad
\dot x=-A\omega\sin\omega t+B\omega\cos\omega t.\]
由$x(0)=0,\dot x(0)=0$依次得$A=-F_0/k$、$B=0$，因果阶跃响应便是$(F_0/k)(1-\cos\omega t)H(t)$。按三个开启时刻线性叠加，全时间解为
\[\begin{aligned}
x(t)=\frac{F_0}{k}\{&(1-\cos\omega t)H(t)\\
&-2[1-\cos\omega(t-\tau)]H(t-\tau)\\
&+[1-\cos\omega(t-2\tau)]H(t-2\tau)\}.
\end{aligned}\]
每个新增响应及其一阶导数在自身开启时都为零，保证各接点的$x,\dot x$连续；加速度按外力跃变改变。逐段求解时，这些连续的位置和速度就是下一段的初值。

对$t\ge2\tau$，三个常数项相消。记$\theta=\omega(t-\tau)$、$h=\omega\tau$，余下余弦组合为
\[-\cos(\theta+h)+2\cos\theta-\cos(\theta-h)
=2(1-\cos h)\cos\theta=4\sin^2(h/2)\cos\theta.\]
所以
\[\boxed{x(t)=\frac{4F_0}{k}\sin^2\frac{\omega\tau}{2}
\cos[\omega(t-\tau)]}.\]
于是撤力后的振幅与机械能为
\[\boxed{A=\frac{4F_0}{k}\sin^2\frac{\omega\tau}{2}},\qquad
\boxed{E=\frac12kA^2=\frac{8F_0^2}{k}\sin^4\frac{\omega\tau}{2}}.\]
非零持续时间下，完全没有余振当且仅当$\omega\tau=2\pi n$，$n=1,2,\ldots$。这要求两段各自持续整数个固有周期；外力面积相消本身没有包含这一相位条件。

图~\ref{fig:c4-v2-s}取$\omega\tau=\pi$。在$2\tau$处虽然速度为零，位移却为$-4F_0/k$，随后仍有振幅$4F_0/k$的自由振动。动力学中的总合力还含弹簧力，不能把外推力的零冲量当成质点动量或整个状态复原的充分条件。对很短的脉冲，$A=F_0\tau^2/m+O(\omega^2\tau^4F_0/m)$，余振虽小，仍不恒为零。
\begin{figure}[H]\centering
\begin{tikzpicture}[>=Stealth,line cap=round,line join=round]

\begin{scope}[x=.58cm,y=.42cm]
\draw[axis](0,0)--(13.3,0)node[right]{$\omega t$};\draw[axis](0,-4.8)--(0,4.9)node[above]{$kx/F_0$};
\draw[curve,domain=0:3.14159265,samples=81]plot(\x,{1-cos(\x r)});
\draw[curve,domain=3.14159265:6.2831853,samples=81]plot(\x,{-1-3*cos(\x r)});
\draw[curve,domain=6.2831853:12.8,samples=121]plot(\x,{-4*cos(\x r)});
\draw[guide](3.14159265,-4.5)--(3.14159265,4.5);\draw[guide](6.2831853,-4.5)--(6.2831853,4.5);
\node[above]at(3.14159265,0){$\pi$};\node[above right]at(6.2831853,0){$2\pi$};\fill(6.2831853,-4)circle(1.5pt);
\end{scope}

\end{tikzpicture}
\caption{取$\omega\tau=\pi$；两次外力总冲量为零，撤力时仍停在远离平衡点的位置。}\label{fig:c4-v2-s}
\end{figure}
\end{solution}
